\chapter{Synchronization}
\label{ch:sync}

\begin{chapterintro}
Every detector in the previous two chapters assumed that the receiver knew exactly when each
symbol began, exactly what the carrier phase was, and exactly where the frame started. No real
receiver knows any of these things. Its oscillator runs at a slightly different frequency from the
transmitter's, its sample clock drifts, the propagation delay is unknown and changing, and the
signal may arrive as a burst out of silence. Synchronization is the art of estimating these
nuisance parameters from the signal itself, quickly and accurately enough that the detector never
notices they were unknown. This chapter treats the receiver as a chain of estimators---frequency,
timing, phase and frame---and develops each from estimation theory and the phase-locked loop. We
derive the Cram\'er--Rao bounds that no synchronizer can beat, design loops that come close to them,
and study the algorithms that real modems use: Costas loops and $M$th-power estimators, the Kay,
Fitz and Luise--Reggiannini frequency estimators, Gardner and Mueller--M\"uller timing detectors with
Farrow interpolators, Schmidl--Cox and the 802.11 preamble, the 5G synchronization signal block,
clock and data recovery in SerDes links, and finally the time and frequency distribution that holds
an entire cellular network to within a microsecond and a half.
\end{chapterintro}

\section{The synchronization problem}
\label{sec:ch10:problem}

\subsection{What the receiver does not know}

Write the complex envelope of a linearly modulated signal as it arrives at the receiver's
matched filter:
\begin{equation}
r(t)=e^{j(2\pi\Delta f\,t+\theta)}\sum_k a_k\,g(t-kT-\tau)+w(t),
\label{eq:ch10:model}
\end{equation}
where $a_k$ are the data symbols, $g(t)$ is the transmit pulse, $T$ the symbol period and $w(t)$
complex white Gaussian noise with power spectral density $N_0$. Four parameters stand between this
waveform and the detector of Chapter~\ref{ch:modulation}:
\begin{description}
\item[Carrier frequency offset (CFO) $\Delta f$.] The transmitter's and receiver's local
oscillators are derived from different crystals, and the channel may add a Doppler shift. A
relative error of $\epsilon$ at carrier frequency $f_c$ produces $\Delta f=\epsilon f_c$.
\item[Carrier phase $\theta$.] Even with the frequencies matched, the phase depends on the
path length to within a fraction of a wavelength---a few centimetres at gigahertz frequencies---and
is therefore effectively random.
\item[Symbol timing $\tau$.] The propagation delay is unknown, and the receiver's sample clock
runs at a slightly different rate from the transmitter's symbol clock, so $\tau$ is not even
constant: it drifts as $\tau(t)=\tau_0+\epsilon_s t$ for a sample clock offset $\epsilon_s$.
\item[Frame timing.] The receiver must also find where the packet, frame or codeword begins in the
stream of symbols, so that bits are handed to the decoder in the right order and the right
bits are interpreted as headers.
\end{description}
The amplitude (the channel gain) is a fifth unknown, handled by automatic gain control and channel
estimation, and in multipath channels the ``phase'' becomes a whole channel impulse response; those
are the subjects of Chapters~\ref{ch:channels} and~\ref{ch:equalization}. Here we assume a flat
channel and concentrate on time and frequency.

Figure~\ref{fig:ch10_impairments} shows what each impairment does to a QPSK constellation sampled
at the output of the matched filter. A static phase offset rotates the constellation; a frequency
offset makes it spin, smearing the points into a ring; a timing offset leaves the points in their
quadrants but spreads each one into a cloud of intersymbol interference (Chapter~\ref{ch:baseband});
and a clock offset makes the timing error sweep slowly through the symbol, so the cloud breathes
open and closed. The detector sees all four at once.

\mdcfig{ch10_impairments}{A QPSK constellation (2000 symbols, $E_s/N_0=22$\,dB, RRC $\beta=0.35$)
under each synchronization impairment. A frequency offset of only $4\times10^{-4}$ of the symbol
rate turns the constellation into a ring within 2000 symbols.}

\subsection{How large are the offsets?}

It is worth knowing the numbers, because they determine which algorithms are feasible.
Table~\ref{tab:ch10:offsets} collects typical values. Two lessons stand out. First, frequency
offsets are often a large fraction of the symbol rate, especially for narrowband signals at high
carrier frequencies: a 2\,ppm crystal at 3.5\,GHz is 7\,kHz, which would rotate a 10\,ksym/s signal by
$250^\circ$ per symbol. Second, the ratio of the offset to the symbol rate is what matters to the
algorithms, so the same crystal is harmless for a 30\,Msym/s satellite carrier and crippling for a
narrowband IoT link.

\begin{table}[htbp]\centering\small
\caption{Representative synchronization offsets. Values are approximate and are meant to show orders
of magnitude; consult the relevant standard for requirements.}
\label{tab:ch10:offsets}
\begin{tabular}{>{\raggedright\arraybackslash}p{4.1cm}>{\raggedright\arraybackslash}p{4.3cm}>{\raggedright\arraybackslash}p{4.9cm}}
\toprule
System & Source of offset & Typical magnitude \\
\midrule
5G NR phone at 3.5\,GHz, before lock & TCXO, $\approx\pm2$\,ppm & $\approx\pm7$\,kHz ($\approx$ 23\% of a 30\,kHz subcarrier)\\
5G NR base station & TS 38.104 requirement, $\pm0.05$\,ppm & $\pm175$\,Hz at 3.5\,GHz \\
Wi-Fi (802.11a/g/n/ac) & $\pm20$\,ppm per device & up to $\pm232$\,kHz at 5.8\,GHz \\
DVB-S2 home terminal & LNB oscillator drift & up to a few MHz \\
LEO satellite, Ku band & Doppler, 7.5\,km/s orbital speed & several hundred kHz \\
GPS L1 receiver, stationary & satellite Doppler & about $\pm5$\,kHz \\
USRP B200 pair & two TCXOs, $\pm2$\,ppm each & up to 4\,ppm: 3.6\,kHz at 900\,MHz \\
\bottomrule
\end{tabular}
\end{table}

\subsection{The synchronizer chain}

A practical receiver does not estimate all parameters jointly---that would be optimal but
prohibitively complex. It estimates them in a sequence, each stage cleaning up the signal enough
for the next to work. Figure~\ref{fig:ch10_chain} shows the typical order for a single-carrier modem.
\begin{enumerate}
\item \textbf{Automatic gain control} brings the signal into the linear range of the ADC and gives
the later stages a known amplitude.
\item \textbf{Coarse frequency correction} comes first, because a large CFO smears the signal
across the matched filter's passband and destroys the operation of everything downstream. Coarse
estimators must have a wide acquisition range and need not be very accurate: a
frequency-locked loop or a non-data-aided estimator is typical.
\item \textbf{Matched filtering and symbol timing recovery} follow. The best timing detectors are
insensitive to carrier phase, so timing can be recovered before the phase is known.
\item \textbf{Frame synchronization} locates a known preamble or unique word, which also provides
data-aided estimates of fine frequency and phase and resolves the phase ambiguity of blind
carrier loops.
\item \textbf{Fine frequency and phase tracking}, usually a decision-directed or pilot-aided
phase-locked loop, removes the residual offset and follows oscillator phase noise for the rest of
the frame.
\end{enumerate}
The ordering is a design choice, not a law. OFDM receivers (Chapter~\ref{ch:ofdm}) find the frame
first, with a preamble designed to tolerate large frequency offsets, and then estimate the CFO from
the same preamble. Burst modems that store a whole packet in memory can iterate, running the chain
twice. But the principle is universal: \emph{each estimator must tolerate the residual errors of
the ones before it.}

\begin{figure}[htbp]\centering
\begin{tikzpicture}[node distance=0.35cm,
  blk/.style={draw=navy,fill=navy!6,thick,minimum height=1.05cm,minimum width=1.45cm,align=center,font=\sffamily\scriptsize},
  est/.style={draw=accent,fill=accent!6,thick,minimum height=0.8cm,minimum width=1.45cm,align=center,font=\sffamily\scriptsize},
  arr/.style={-{Stealth[length=2mm]},thick,navy}]
\node[blk] (adc) {ADC\\+ AGC};
\node[blk,right=of adc] (cfo) {coarse\\CFO\\correction};
\node[blk,right=of cfo] (mf) {matched\\filter};
\node[blk,right=of mf] (int) {timing\\interpolator};
\node[blk,right=of int] (rot) {phase\\rotator};
\node[blk,right=of rot] (det) {detector\\+ decoder};
\node[est,below=0.55cm of cfo] (fll) {FLL or\\feedforward};
\node[est,below=0.55cm of int] (ted) {TED + loop\\filter + NCO};
\node[est,below=0.55cm of rot] (pll) {PLL / pilots};
\node[est,above=0.45cm of rot] (fs) {frame sync\\(preamble)};
\draw[arr] (adc)--(cfo); \draw[arr] (cfo)--(mf); \draw[arr] (mf)--(int); \draw[arr] (int)--(rot); \draw[arr] (rot)--(det);
\draw[arr,accent] (cfo.south)++(0.25,0) -- ++(0,-0.55);
\draw[arr,accent] (fll.north)++(-0.25,0) -- ++(0,0.55);
\draw[arr,accent] (int.south)++(0.25,0) -- ++(0,-0.55);
\draw[arr,accent] (ted.north)++(-0.25,0) -- ++(0,0.55);
\draw[arr,accent] (rot.south)++(0.25,0) -- ++(0,-0.55);
\draw[arr,accent] (pll.north)++(-0.25,0) -- ++(0,0.55);
\draw[arr,accent] ($(int.east)!0.5!(rot.west)$) |- (fs.west);
\draw[arr,accent] (fs.east) -| ($(det.north)+(-0.3,0)$);
\draw[arr,accent] (fs.south) -- (rot.north);
\end{tikzpicture}
\caption{The synchronizer chain of a single-carrier receiver. Signal path in navy; estimators in red.
Each estimator operates on a signal already partly corrected by the stages before it. The frame
synchronizer supplies the data-aided phase reference that resolves the carrier loop's ambiguity and
tells the decoder where the codewords start.}
\label{fig:ch10_chain}
\end{figure}

\subsection{A vocabulary of synchronizers}

Three distinctions recur throughout the chapter.

\paragraph{Feedback or feedforward.} A \term{feedback} synchronizer (a loop) measures the residual
error \emph{after} correction and nudges its estimate to reduce it; the phase-locked loop is the
archetype. A \term{feedforward} synchronizer computes the estimate directly from a block of received
samples and applies it. Loops are cheap, track slow variations naturally and are ideal for
continuous transmission, but they take time to acquire and can hang or slip. Feedforward estimators
acquire in one block, which suits bursts, but must be re-run to track drift.

\paragraph{Data-aided, decision-directed or blind.} A \term{data-aided} (DA) estimator knows the
transmitted symbols---a preamble or pilots---and so can strip the modulation exactly. A
\term{decision-directed} (DD) estimator substitutes the detector's decisions for the true symbols;
it performs like a DA estimator when decisions are reliable and fails catastrophically when they are
not. A \term{non-data-aided} (NDA), or blind, estimator removes the modulation by a nonlinearity that
is insensitive to the data, such as raising PSK symbols to the $M$th power. It needs no overhead,
but pays with worse accuracy at low SNR and with ambiguities.

\paragraph{Acquisition and tracking.} Getting from total ignorance to a small error is
\emph{acquisition}; keeping the error small as parameters drift is \emph{tracking}. They pull in
opposite directions: acquisition wants wide bandwidth and wide range, tracking wants narrow bandwidth
to suppress noise. Most good designs therefore use two modes or two estimators: a coarse,
wide-range one and a fine, narrow-band one, or a loop whose bandwidth is reduced after lock
(\emph{gear shifting}).

\begin{history}[De Bellescize and the synchronous receiver, 1932]
The phase-locked loop is older than most of the electronics it now lives in. In 1932 the French
engineer Henri de Bellescize published ``La r\'eception synchrone'' in \emph{L'Onde \'Electrique},
describing a receiver whose local oscillator was locked in phase to the incoming carrier by
feedback, so that AM could be demodulated coherently---a homodyne receiver made practical. The idea
was ahead of its components: vacuum-tube oscillators drifted, and the superheterodyne with its
envelope detector was cheaper and forgiving. The PLL spent the next two decades in niches such as
the horizontal-sync ``flywheel'' circuits of television receivers, which kept the scanning beam
steady through noise bursts. It came of age in the 1950s, when colour television and the space
programme gave it two jobs that nothing else could do, which we shall meet shortly.
\end{history}

\section{Estimation theory for synchronizers}
\label{sec:ch10:estimation}

Synchronization is parameter estimation. Before designing any algorithm it pays to ask two
questions: what is the best possible estimator, and how good can any estimator be? Maximum
likelihood answers the first, the Cram\'er--Rao bound the second.

\subsection{The likelihood function}

Sample the output of the matched filter at candidate timing $\tau$ after de-rotating by a candidate
frequency $\nu$, giving $z_k(\tau,\nu)$. For a known block of $N$ symbols $\vect{a}$ in complex AWGN,
the log-likelihood of the parameters $(\theta,\nu,\tau)$ reduces---after dropping terms that do not
depend on them---to a correlation between the received samples and the known symbols:
\begin{equation}
\Lambda(\theta,\nu,\tau)=\frac{2}{N_0}\,\re\Bigl\{e^{-j\theta}\sum_{k=0}^{N-1}a_k^*\,z_k(\tau,\nu)\Bigr\}.
\label{eq:ch10:llf}
\end{equation}
Everything in this chapter is an approximation to maximising~\eqref{eq:ch10:llf}. Maximising over
$\theta$ alone is easy: the real part of $e^{-j\theta}X$ is largest when $\theta=\arg X$, so the
\term{maximum-likelihood phase estimate} is
\begin{equation}
\hat\theta=\arg\sum_{k=0}^{N-1}a_k^*\,z_k .
\label{eq:ch10:mlphase}
\end{equation}
Substituting it back, the likelihood for frequency and timing becomes
$\bigl|\sum_k a_k^*z_k(\tau,\nu)\bigr|$, the magnitude of a correlation. Maximising over $\nu$ is a
search for the peak of a Fourier transform of $a_k^*z_k$; maximising over $\tau$ is a search for
the peak of a correlation in time. Frame synchronization, frequency estimation, timing estimation
and phase estimation are, at bottom, the same operation: \emph{correlate with what you know and find
the peak.}

When the data are unknown, the likelihood must be averaged over all possible symbol sequences. For
BPSK the average of $e^{\re\{x\}}$ and $e^{-\re\{x\}}$ is $\cosh(\re\{x\})$, and for low SNR
$\ln\cosh x\approx x^2/2$: the blind likelihood becomes a sum of \emph{squares} of the samples, which
is why squaring and $M$th-power nonlinearities appear in every NDA synchronizer. At high SNR,
$\ln\cosh x\approx|x|$, which is decision-direction in disguise.

\subsection{The Cram\'er--Rao bound}

For any unbiased estimator $\hat\lambda$ of a parameter $\lambda$ from observations $\vect{r}$, the
\term{Cram\'er--Rao bound} (CRB) states
\begin{equation}
\operatorname{var}(\hat\lambda)\ \ge\ \frac{1}{J(\lambda)},\qquad
J(\lambda)=-\E\Bigl[\frac{\partial^2\ln p(\vect r;\lambda)}{\partial\lambda^2}\Bigr],
\end{equation}
where $J$ is the \emph{Fisher information}: the expected curvature of the log-likelihood. A sharply
peaked likelihood carries a lot of information about where its peak is. For the model
\eqref{eq:ch10:model} with known data, the calculation is a few lines of calculus, and the results
are the three bounds every synchronizer designer carries in their head
(Table~\ref{tab:ch10:crb}).

\begin{table}[htbp]\centering\small
\caption{Cram\'er--Rao bounds for data-aided synchronization with $N$ (or $L_0$) symbols at
symbol-energy-to-noise ratio $\esno$. The frequency $\nu=\Delta fT$ is in cycles per symbol and
$\xi$ is the normalised mean-square bandwidth of the pulse.}
\label{tab:ch10:crb}
\begin{tabular}{lll}
\toprule
Parameter & Bound & Scaling \\
\midrule
Phase (frequency known) & $\operatorname{var}(\hat\theta)\ge\dfrac{1}{2N\,\esno}$ & $1/N$ \\[2ex]
Frequency (phase unknown) & $\operatorname{var}(\hat\nu)\ge\dfrac{3}{2\pi^2N(N^2-1)\,\esno}$ & $1/N^3$ \\[2ex]
Timing & $\operatorname{var}(\hat\tau/T)\ge\dfrac{1}{8\pi^2\xi\,L_0\,\esno}$ & $1/L_0$ \\[1.5ex]
\bottomrule
\end{tabular}
\end{table}

The phase bound is intuitive: averaging $N$ samples, each with phase noise variance
$1/(2\esno)$, reduces the variance by $N$. The frequency bound is more interesting: it improves as
$N^3$, not $N$. Frequency is the slope of phase against time, and fitting a slope to a longer
record improves in two ways at once: there are more points, and they are spread further apart. The
lever arm matters. Doubling the preamble length improves frequency accuracy by a factor of
$2^{3/2}\approx2.8$ in standard deviation.

For timing, the parameter
\begin{equation}
\xi=T^2\,\frac{\int f^2|G(f)|^2\,df}{\int|G(f)|^2\,df}
\end{equation}
measures how much high-frequency content the pulse has. A pulse with sharp edges locates itself
precisely in time; a smooth one does not. For a root-raised-cosine pulse of roll-off $\beta$,
$\xi=\tfrac1{12}+\beta^2\bigl(\tfrac14-\tfrac{2}{\pi^2}\bigr)$, which grows only from $0.083$ to $0.13$
as $\beta$ goes from 0 to 1. This is the timing-recovery analogue of the uncertainty principle:
bandwidth buys timing accuracy, and it is the same principle that makes wideband GNSS signals and
radar chirps precise.

\subsection{The modified Cram\'er--Rao bound}

When the data are random and unknown the true CRB is hard to compute, because the likelihood is
an average over all sequences. D'Andrea, Mengali and Reggiannini (1994) introduced the
\term{modified Cram\'er--Rao bound} (MCRB), in which the expectation over data is taken
\emph{outside} the Fisher information. The result is simple: the MCRB has exactly the formulas of
Table~\ref{tab:ch10:crb}, as if the data were known. The MCRB is looser than (below) the true CRB,
so it is always a valid lower bound; it equals the CRB for data-aided estimation, and blind or
decision-directed estimators approach it only at high SNR. In practice the MCRB is the yardstick
against which almost every synchronizer in the literature is plotted, including those in this
chapter.

\begin{worked}[How long must a preamble be?]
A burst modem sends QPSK at $R_s=1$\,Msym/s and must estimate its CFO from a 64-symbol preamble at
$\esno=10$\,dB. The frequency CRB is
\[
\sigma_\nu\ge\sqrt{\frac{3}{2\pi^2\cdot64\cdot(64^2-1)\cdot10}}\approx2.4\times10^{-4}\ \text{cycles/symbol},
\]
or 240\,Hz. That seems tiny, but the phase error it causes grows linearly: after 1000 data symbols
the phase has drifted by $2\pi\times2.4\times10^{-4}\times1000\approx1.5$\,rad (one standard
deviation), far beyond what QPSK tolerates. Even a perfect preamble estimator must therefore be
followed by a phase tracker, or by pilots inserted through the payload, or by a longer preamble.
Doubling the preamble to 128 symbols cuts $\sigma_\nu$ by $2.8$, still not enough without tracking
for long packets. This is the fundamental reason that almost every burst modem combines a
feedforward preamble estimate with a tracking loop or pilot-aided phase correction.
\end{worked}

\subsection{Thresholds and outliers}

The CRB describes small errors. At low SNR every nonlinear estimator exhibits a \term{threshold}:
below some SNR, noise occasionally produces a correlation peak larger than the true one, and the
estimate jumps to a completely wrong value---an \emph{outlier}. Outliers are rare but enormous, so
the mean-square error rises abruptly above the bound. The threshold, not the bound, usually decides
whether a design works. Figure~\ref{fig:ch10_freq_est} later in the chapter shows thresholds for four
frequency estimators with the same asymptotic accuracy but wildly different thresholds. The design
rule is: \emph{choose an algorithm whose threshold lies comfortably below the operating SNR of the
code}. With modern LDPC and turbo codes operating at $\esno$ near or below 0\,dB, this is often the
hardest constraint in the whole receiver.

\begin{pitfall}[Synchronizing below the code's threshold]
A rate-1/4 LDPC code in DVB-S2 decodes QPSK at $\esno\approx-2.4$\,dB. The coding gain is
worthless if the synchronizers fail at 0\,dB. Decision-directed loops are useless there, because
two symbols in ten are wrong before decoding, and blind $M$th-power estimators lose heavily to noise.
This is exactly why DVB-S2 inserts known pilot blocks and a $\pi/2$-BPSK frame header: the
synchronizers need known symbols to work at SNRs where the data symbols are mostly noise. Always
check that synchronization works at the lowest SNR your code does, and plot its performance with
real coded BER, not against an ideal-sync curve.
\end{pitfall}

\section{The phase-locked loop}
\label{sec:ch10:pll}

The phase-locked loop is the single most important circuit in synchronization. It recovers
carriers, recovers clocks, synthesises every local oscillator in every radio, demodulates FM,
and appears, in digital disguise, in the timing and carrier loops of every modem. Its theory
repays careful study, because the same equations describe every feedback synchronizer in this
chapter.

\subsection{Anatomy and the nonlinear model}

A PLL has three parts (Figure~\ref{fig:ch10_pll_model}). A \term{phase detector} (PD) compares the
phase $\theta(t)$ of the input with the phase $\hat\theta(t)$ of a local oscillator and produces an
error signal $e(t)=K_d\,S(\theta-\hat\theta)$, where $S(\cdot)$ is the detector's characteristic or
\term{S-curve}. A \term{loop filter} $F(s)$ smooths the error. A \term{voltage-controlled
oscillator} (VCO), or in digital form a numerically controlled oscillator, changes its frequency in
proportion to the filtered error, so its phase is the integral: $d\hat\theta/dt=K_0\,v(t)$.

The simplest phase detector is a multiplier. Multiplying $\cos(\omega t+\theta)$ by
$-\sin(\omega t+\hat\theta)$ and low-pass filtering leaves $\tfrac12\sin(\theta-\hat\theta)$. In
complex baseband the equivalent operation is $\im\{z\,e^{-j\hat\theta}\}$ for $z=e^{j\theta}$, and the
S-curve is $S(\phi)=\sin\phi$ with phase error $\phi=\theta-\hat\theta$. The loop is therefore
described by the nonlinear differential equation
\begin{equation}
\frac{d\hat\theta}{dt}=K_0K_d\,\bigl[f*\sin(\theta-\hat\theta)\bigr](t),
\end{equation}
where $f(t)$ is the impulse response of the loop filter. It has no closed-form solution in general but
is easy to simulate. Near lock, $\sin\phi\approx\phi$ and the loop becomes linear.

\begin{figure}[htbp]\centering
\begin{tikzpicture}[node distance=0.9cm,
  blk/.style={draw=navy,fill=navy!6,thick,minimum height=0.95cm,minimum width=1.6cm,align=center,font=\sffamily\small},
  sum/.style={draw=navy,thick,circle,inner sep=1.5pt,font=\small},
  arr/.style={-{Stealth[length=2.2mm]},thick,navy}]
\node (in) {$\theta(t)$};
\node[sum,right=0.7cm of in] (s) {$\Sigma$};
\node[blk,right=of s] (pd) {$K_d$\\[-1pt]\scriptsize phase detector};
\node[blk,right=of pd] (lf) {$F(s)$\\[-1pt]\scriptsize loop filter};
\node[blk,right=of lf] (vco) {$K_0/s$\\[-1pt]\scriptsize VCO / NCO};
\node[right=0.8cm of vco] (out) {$\hat\theta(t)$};
\draw[arr] (in)--node[above,pos=0.8,font=\scriptsize]{$+$}(s);
\draw[arr] (s)--node[above,font=\scriptsize]{$\phi$}(pd);
\draw[arr] (pd)--node[above,font=\scriptsize]{$e$}(lf);
\draw[arr] (lf)--node[above,font=\scriptsize]{$v$}(vco);
\draw[arr] (vco)--(out);
\draw[arr] ($(vco.east)+(0.4,0)$) -- ++(0,-1.1) -| node[left,pos=0.9,font=\scriptsize]{$-$} (s.south);
\end{tikzpicture}
\caption{Linearised model of a phase-locked loop. The phase detector is replaced by its slope
$K_d$ at the lock point; the oscillator integrates its control signal.}
\label{fig:ch10_pll_model}
\end{figure}

\subsection{The linear model}

With $\sin\phi\approx\phi$, the loop in Figure~\ref{fig:ch10_pll_model} is an ordinary linear feedback
system. Its open-loop gain is $G(s)=K_0K_dF(s)/s$, and the transfer functions from input phase to
output phase and to phase error are
\begin{equation}
H(s)=\frac{\hat\Theta(s)}{\Theta(s)}=\frac{G(s)}{1+G(s)},\qquad
\frac{\Phi(s)}{\Theta(s)}=1-H(s)=\frac{1}{1+G(s)}.
\label{eq:ch10:H}
\end{equation}
$H(s)$ is low-pass: the loop follows slow changes in the input phase and ignores fast ones. The
error transfer function $1-H(s)$ is high-pass: it passes fast input phase fluctuations straight
into the error. These two facts summarise the loop's two jobs and their conflict. Noise on the input
is fast, and we want the loop to ignore it (narrow $H$); oscillator phase noise and Doppler are
slow, and we want the loop to follow them (wide $H$).

\subsection{Order and type}

The \term{order} of the loop is the order of the denominator of $H(s)$; its \term{type} is the number
of perfect integrators in $G(s)$. The VCO always contributes one. With $F(s)=1$ (no filter) the
loop is first order and type 1; with a proportional-plus-integral filter,
\begin{equation}
F(s)=\frac{1+s\tau_2}{s\tau_1},
\end{equation}
it is second order and type 2, the workhorse of carrier recovery. Writing $K=K_0K_d$,
\begin{equation}
H(s)=\frac{2\zeta\omega_ns+\omega_n^2}{s^2+2\zeta\omega_ns+\omega_n^2},\qquad
\omega_n=\sqrt{K/\tau_1},\quad \zeta=\tfrac12\,\omega_n\tau_2,
\label{eq:ch10:H2}
\end{equation}
with \term{natural frequency} $\omega_n$ and \term{damping factor} $\zeta$. Type matters because it
determines which inputs the loop can track with zero steady-state error. The final-value theorem
applied to $\Phi(s)=[1-H(s)]\Theta(s)$ gives Table~\ref{tab:ch10:steady}.

\begin{table}[htbp]\centering\small
\caption{Steady-state phase error of linear PLLs. A frequency step $\Delta\omega$ is a phase ramp; a
frequency ramp $\dot\omega$ (for example the Doppler rate of a satellite pass) is a parabola in phase.}
\label{tab:ch10:steady}
\begin{tabular}{llll}
\toprule
Input & Type 1 (first order) & Type 2 (PI filter) & Type 3 \\
\midrule
Phase step $\Delta\theta$ & 0 & 0 & 0 \\
Frequency step $\Delta\omega$ & $\Delta\omega/K$ & 0 & 0 \\
Frequency ramp $\dot\omega$ & grows without bound & $\dot\omega/\omega_n^2$ & 0 \\
\bottomrule
\end{tabular}
\end{table}

A type-1 loop tracks a frequency offset only with a standing phase error $\Delta\omega/K$; for a
coherent demodulator that error is a permanent rotation of the constellation. A type-2 loop's
integrator absorbs the frequency offset completely: the integrator's value \emph{is} the frequency
estimate. A type-2 loop still lags a frequency ramp, and receivers for low-orbit satellites or
high-dynamics GNSS, where the Doppler rate is large, use type-3 loops or aid the loop with a
predicted Doppler.

Figure~\ref{fig:ch10_pll_response} shows the transient response of the type-2 loop. For a phase step
the error decays with an overshoot that depends on $\zeta$; for a frequency step the error rises and
then returns to zero as the integrator charges up. The value $\zeta=1/\sqrt2\approx0.707$ is the
customary compromise between speed and overshoot, and it also gives a loop with a gently peaked
$|H|$ that passes noise without much amplification. Lightly damped loops ($\zeta=0.3$) ring and
amplify noise near $\omega_n$ (the peak in panel (c)); heavily damped ones approach first-order
behaviour and are slow to remove frequency errors.

\mdcfig{ch10_pll_response}{The second-order type-2 PLL. (a) Phase error after a unit phase step.
(b) Phase error after a frequency step, normalised by $\Delta\omega/\omega_n$: the type-2 loop returns
to zero, while a type-1 loop with gain $K=\omega_n$ settles at a standing error. (c) Closed-loop
$|H(j\omega)|$ and error response $|1-H(j\omega)|$.}

\subsection{Noise bandwidth and phase jitter}

Suppose the input is a carrier of power $C$ in white noise of density $N_0$. The noise at the phase
detector output is equivalent to a phase noise of two-sided density $N_0/(2C)$\,rad$^2$/Hz, and the
loop passes it through $H$. The variance of the output phase is therefore
\begin{equation}
\sigma_\phi^2=\frac{N_0}{2C}\int_{-\infty}^{\infty}|H(j2\pi f)|^2\,df=\frac{N_0B_L}{C}=\frac{1}{\rho},
\qquad B_L=\int_0^\infty|H(j2\pi f)|^2\,df,
\label{eq:ch10:jitter}
\end{equation}
where $B_L$ is the one-sided \term{loop noise bandwidth} and $\rho=C/(N_0B_L)$ is the \term{loop SNR}:
the SNR of the carrier measured in the loop bandwidth. For the second-order loop,
\begin{equation}
B_L=\frac{\omega_n}{2}\Bigl(\zeta+\frac{1}{4\zeta}\Bigr),
\label{eq:ch10:BL}
\end{equation}
which is $0.53\,\omega_n$ at $\zeta=0.707$. In a digital loop that updates once per symbol, with a
data-aided phase detector, the same result reads
\begin{equation}
\sigma_\phi^2=\frac{B_LT}{\esno},
\label{eq:ch10:jitter_digital}
\end{equation}
where $B_LT$ is the noise bandwidth normalised to the symbol rate. Typical carrier loops use
$B_LT$ between $10^{-3}$ and $10^{-2}$: at $\esno=10$\,dB and $B_LT=0.005$, $\sigma_\phi=0.022$\,rad,
or $1.3^\circ$. Figure~\ref{fig:ch10_jitter}(a) confirms~\eqref{eq:ch10:jitter_digital} by simulation
with the \texttt{commlib} loop. The decision-directed loop matches the data-aided one at high SNR
and departs from it below about 8\,dB, where decision errors inject additional noise.

Equation~\eqref{eq:ch10:jitter_digital} also says something about optimality: a loop of noise
bandwidth $B_L$ behaves like an averager over roughly $1/(2B_LT)$ symbols, and the phase CRB for that
many symbols is $B_LT/\esno$. A well-designed loop is an efficient estimator; the only freedom is
how long a memory to give it.

\mdcfig{ch10_jitter}{(a) Phase-error variance of a second-order QPSK loop against $\esno$ for three
loop bandwidths: lines are $B_LT/(\esno)$, circles a data-aided simulation, crosses the
decision-directed \texttt{commlib.pll\_dd}. (b) At $\esno=3$\,dB the decision-directed loop
cycle-slips: its phase estimate wanders by whole multiples of $\pi/2$.}

\begin{worked}[Designing a carrier loop]
A telemetry receiver must track a residual carrier with $C/N_0=40$\,dB-Hz. Choose
$B_L=50$\,Hz and $\zeta=0.707$. From~\eqref{eq:ch10:BL}, $\omega_n=2B_L/(\zeta+1/4\zeta)=94.3$\,rad/s
($f_n=15$\,Hz). The loop SNR is $\rho=10^4/50=200$ (23\,dB), so $\sigma_\phi=1/\sqrt{200}=0.071$\,rad,
or $4^\circ$---comfortable for BPSK telemetry. A Doppler rate of $\dot f=20$\,Hz/s gives a
steady-state error $\dot\omega/\omega_n^2=2\pi\cdot20/94.3^2=0.014$\,rad, negligible. But if the
initial frequency uncertainty is $\pm1$\,kHz, the loop will take a very long time to pull in, as we
now see; the receiver needs a frequency acquisition aid.
\end{worked}

\subsection{Acquisition: lock-in, pull-in and hang-up}

The linear model says nothing about how the loop gets into lock. For that we need the nonlinear
model, and the classical results (Gardner's \emph{Phaselock Techniques} is the standard source) are
these:
\begin{itemize}
\item The \term{lock-in range} is the frequency offset from which the loop locks within about one
beat cycle, without slipping. For a second-order loop it is approximately $\Delta\omega_L\approx2\zeta\omega_n$.
\item Beyond lock-in, a type-2 loop with a perfect integrator still \emph{pulls in} from any offset
within the VCO's range, because the beat note's slight asymmetry slowly charges the integrator. But
pull-in is slow: the time grows as the square of the offset and the inverse cube of the natural
frequency, approximately
\begin{equation}
T_p\approx\frac{\Delta\omega^2}{2\zeta\omega_n^3}
\label{eq:ch10:pullin}
\end{equation}
for a sinusoidal phase detector. During pull-in the loop slips cycles repeatedly.
\item The \term{hold range} is the offset beyond which a locked loop loses lock when the input
frequency changes slowly. For the ideal type-2 loop it is limited only by the VCO and the
integrator's dynamic range.
\item A \term{hang-up} occurs when the loop starts near the unstable equilibrium of the S-curve
($\phi=\pi$ for a sinusoidal detector), where the restoring force is nearly zero. The loop can
linger there for a long time before noise pushes it off; a sawtooth (arctangent) detector, whose
restoring force is largest near $\pm\pi$, largely cures it.
\end{itemize}
For the loop of the last worked example, a 1\,kHz offset gives
$T_p\approx(2\pi\cdot1000)^2/(2\cdot0.707\cdot94.3^3)\approx33$\,s, which is useless. Practical receivers
therefore acquire frequency by other means---a frequency sweep of the NCO, a frequency-locked loop,
or a feedforward FFT estimate (Section~\ref{sec:ch10:freq})---and hand the PLL an offset well inside its
lock-in range. Figure~\ref{fig:ch10_acquisition} shows pull-in of a QPSK decision-directed loop with
an offset of 0.008 cycles/symbol: the wider loops lock almost at once, while the narrow loop slips
fifteen quarter-cycles over several hundred symbols before its integrator reaches the right
frequency.

\mdcfig{ch10_acquisition}{Pull-in of a second-order decision-directed QPSK loop ($\esno=20$\,dB,
offset $\nu=0.008$ cycles/symbol). (a) Phase error in units of the $\pi/2$ ambiguity: the narrowest
loop slips fifteen times before locking. (b) The loop's frequency estimate (NCO increment per symbol)
converging to the true offset.}

\subsection{Cycle slips}

Even in lock, noise occasionally pushes the phase error over the top of the S-curve, and the loop
settles one full cycle (or, for an $M$-fold ambiguous detector, $2\pi/M$) away. The result is a
\term{cycle slip}: in a coherent receiver every subsequent symbol is rotated, and without
differential encoding or a fresh reference every bit is wrong until the slip is detected. Viterbi
solved the first-order loop exactly in 1963 using the Fokker--Planck equation; the mean time between
slips is
\begin{equation}
\bar T_{\text{slip}}=\frac{\pi^2\rho\,I_0^2(\rho)}{2B_L},
\label{eq:ch10:slip}
\end{equation}
where $I_0$ is the modified Bessel function. Because $I_0(\rho)$ grows like $e^\rho/\sqrt{2\pi\rho}$,
the slip rate falls exponentially with loop SNR. The consequence is a sharp threshold.

\begin{worked}[The cliff edge of cycle slipping]
A loop with $B_L=100$\,Hz operates at a loop SNR of $\rho=10$ (10\,dB). Then
$I_0(10)\approx2816$ and $\bar T_{\text{slip}}\approx\pi^2\cdot10\cdot2816^2/200\approx3.9\times10^6$\,s:
about one slip every 45 days. At $\rho=5$ (7\,dB), $I_0(5)\approx27.2$ and
$\bar T_{\text{slip}}\approx180$\,s: one slip every three minutes. A 3\,dB loss of carrier power turns a
loop that almost never slips into one that slips constantly. For an $M$-PSK carrier loop the relevant
SNR is roughly that of the $M$-fold phase, $\rho/M^2$, reduced further by the squaring loss discussed
below, which is why QPSK and 8PSK loops slip at much higher SNR than a pure-carrier loop of the same
bandwidth.
\end{worked}

\subsection{The digital PLL}

In a software or FPGA receiver the PLL runs once per sample or once per symbol, and each analog
component has a digital twin (Figure~\ref{fig:ch10_dpll}). The phase detector is an arithmetic
function of the de-rotated sample---$\im\{z\,a^*\}$ or $\arg(z\,a^*)$. The loop filter is a
proportional-plus-integral update,
\begin{equation}
I[n]=I[n-1]+k_i\,e[n],\qquad v[n]=k_p\,e[n]+I[n],
\end{equation}
and the oscillator is the NCO of Chapter~\ref{ch:dsp}: a phase accumulator
$\hat\theta[n+1]=\hat\theta[n]+v[n]$ addressing a sine table or a CORDIC that rotates the next
sample by $e^{-j\hat\theta}$. The integrator $I[n]$ holds the frequency estimate in radians per
update.

\begin{figure}[htbp]\centering
\begin{tikzpicture}[
  blk/.style={draw=navy,fill=navy!6,thick,minimum height=0.85cm,minimum width=1.35cm,align=center,font=\sffamily\scriptsize},
  mul/.style={draw=navy,thick,circle,inner sep=0.5pt,font=\small},
  arr/.style={-{Stealth[length=2mm]},thick,navy}]
\node (in) at (0,0) {\small $y[n]$};
\node[mul] (m) at (1.2,0) {$\times$};
\node[blk] (pd) at (3.0,0) {phase\\detector};
\node[blk] (kp) at (5.4,0) {$k_p$};
\node[blk] (ki) at (5.4,-1.2) {$k_i/(1-z^{-1})$};
\node[mul] (s) at (7.1,0) {$+$};
\node[blk] (acc) at (8.9,0) {NCO\\$z^{-1}/(1-z^{-1})$};
\node[blk] (exp) at (1.2,-2.3) {$e^{-j\hat\theta}$\\(LUT/CORDIC)};
\node[font=\scriptsize,anchor=west] (out) at (2.3,1.0) {$z[n]$ to detector};
\draw[arr] (in)--(m); \draw[arr] (m)--(pd);
\draw[arr] (pd)--node[above,font=\scriptsize]{$e[n]$}(kp);
\draw[arr] (4.3,0) |- (ki.west);
\draw[arr] (kp)--(s); \draw[arr] (ki.east) -| (s.south);
\draw[arr] (s)--node[above,font=\scriptsize]{$v$}(acc);
\draw[arr] (acc.south) |- node[pos=0.75,below,font=\scriptsize]{$\hat\theta[n]$} (exp.east);
\draw[arr] (exp.north)--(m.south);
\draw[arr] (1.9,0) |- (out.west);
\end{tikzpicture}
\caption{A digital PLL. The loop filter's integrator holds the frequency estimate; the NCO
accumulates phase and rotates the next input sample.}
\label{fig:ch10_dpll}
\end{figure}

For small $B_LT$ the digital gains relate to the analog parameters simply:
$k_p\approx2\zeta\omega_nT$ and $k_i\approx(\omega_nT)^2$, divided by the detector gain $K_d$. The
exact mapping, used by \texttt{commlib} and by Rice's textbook, is
\begin{lstlisting}
def loop_gains(bn, zeta=0.7071, kd=1.0, k0=1.0):
    theta = bn / (zeta + 1 / (4 * zeta))     # = omega_n T / 2
    d = 1 + 2 * zeta * theta + theta ** 2
    kp = 4 * zeta * theta / d / (kd * k0)
    ki = 4 * theta ** 2 / d / (kd * k0)
    return kp, ki
\end{lstlisting}
With $B_LT=0.01$ this gives $k_p=0.0263$ and $k_i=3.5\times10^{-4}$. Two digital effects deserve
attention. First, \emph{loop delay}: pipelined FPGA implementations often apply the correction
several samples after the error is measured, which adds phase lag and erodes the stability margin.
The rule of thumb is to keep the loop bandwidth small compared with $1/(2\pi D)$ for a delay of $D$
updates. Second, the detector gain $K_d$ depends on signal amplitude unless the input is
normalised, so a loop designed for one AGC setting has a different bandwidth at another.

\begin{history}[From the colour burst to the Deep Space Network]
The PLL's breakthrough came from two very different customers in the 1950s. When the US adopted the
NTSC colour television standard in 1953, the chrominance information rode on a suppressed
3.579545\,MHz subcarrier whose phase encoded hue. Each receiver had to regenerate that subcarrier
from a short \emph{colour burst}, about nine cycles transmitted on the back porch of every scan
line, and hold its phase steady for the rest of the line; any phase error showed up on the screen as
a change of colour. The solution was a crystal oscillator phase-locked to the burst, and it put a
PLL in tens of millions of living rooms.

At the same time the Jet Propulsion Laboratory was building tracking receivers for the first
satellites. Jaffe and Rechtin's 1955 paper on phase-lock circuits that perform near-optimally over a
wide range of signal and noise levels laid out the theory, and JPL's ``Microlock'' phase-locked
tracking system followed Explorer~1 in 1958. Their descendants in the Deep Space Network lock onto
spacecraft carriers with loop bandwidths of a few hertz or less. In the Apollo Unified S-Band system
the spacecraft transponder phase-locked to the uplink and retransmitted at exactly 240/221 times its
frequency, so the ground station could measure the two-way Doppler shift coherently and determine
the spacecraft's range rate with great precision.
\end{history}

\subsection{PLL frequency synthesizers}

Chapter~\ref{ch:sdr} introduced the PLL frequency synthesizer, which locks a VCO to $N$ times a
crystal reference. Its modern form, the \term{charge-pump PLL} (Figure~\ref{fig:ch10_cppll}), uses
a \emph{phase-frequency detector} (PFD): two flip-flops clocked by the reference and by the divided
VCO output, whose outputs switch current sources into the loop filter. Unlike a multiplier, the PFD
responds to frequency as well as phase---when the two inputs differ in frequency its average output
has the right sign to pull them together---so a charge-pump PLL acquires from any offset within the
VCO's tuning range without sweeping. The loop filter is usually a series resistor and capacitor
with a small shunt capacitor, giving a type-2 loop with an extra pole for ripple suppression.

The synthesizer's loop bandwidth shapes its phase noise. Inside the loop bandwidth the output follows
the reference, whose phase noise is multiplied by $N$ (raised by $20\log_{10}N$\,dB), together with
the noise of the PFD, charge pump and divider. Outside the bandwidth the loop cannot correct the VCO,
and the VCO's own noise dominates. The optimum bandwidth is roughly where the multiplied reference
noise and the free-running VCO noise cross; fractional-$N$ synthesizers push the loop bandwidth up
by using a higher reference frequency, and must then suppress the sigma-delta modulator's shaped
quantization noise (Chapter~\ref{ch:pcm}) with the same loop filter. The residual phase noise is
what the receiver's carrier recovery loop must then track, and the two loops must be designed
together: a carrier loop narrower than the synthesizer's noise pedestal cannot follow it.

\begin{figure}[htbp]\centering
\begin{tikzpicture}[node distance=0.95cm,
  blk/.style={draw=navy,fill=navy!6,thick,minimum height=0.9cm,minimum width=1.35cm,align=center,font=\sffamily\scriptsize},
  arr/.style={-{Stealth[length=2mm]},thick,navy}]
\node (ref) {\scriptsize $f_{\text{ref}}$};
\node[blk,right=0.5cm of ref] (pfd) {phase-freq.\\detector};
\node[blk,right=of pfd] (cp) {charge\\pump};
\node[blk,right=of cp] (lf) {loop filter\\$R$, $C_1$, $C_2$};
\node[blk,right=of lf] (vco) {VCO};
\node[right=0.5cm of vco,font=\scriptsize,align=left] (out) {$f_{\text{out}}$\\$=Nf_{\text{ref}}$};
\node[blk,below=0.55cm of cp,minimum width=3.5cm] (div) {$\div N$ (or $N+\text{frac}$, driven by $\Sigma\Delta$)};
\draw[arr] (ref)--(pfd); \draw[arr] (pfd)--node[above,font=\scriptsize]{up/dn}(cp);
\draw[arr] (cp)--node[above,font=\scriptsize]{$I_{cp}$}(lf); \draw[arr] (lf)--node[above,font=\scriptsize]{$v$}(vco);
\draw[arr] (vco)--(out);
\draw[arr] ($(vco.east)+(0.25,0)$) |- (div.east);
\draw[arr] (div.west) -| (pfd.south);
\end{tikzpicture}
\caption{A charge-pump PLL frequency synthesizer. The phase-frequency detector makes the loop
self-acquiring; a sigma-delta modulator dithering the divide ratio gives fractional-$N$ resolution.}
\label{fig:ch10_cppll}
\end{figure}

\section{Carrier phase recovery}
\label{sec:ch10:carrier}

With the PLL in hand, carrier recovery becomes a question of which phase detector to use. The
answer depends on what the receiver knows about the symbols and on whether the signal has a carrier
component at all. Modern digital modulations do not: a balanced PSK or QAM constellation has zero
mean, so its spectrum contains no line at the carrier frequency for a PLL to lock onto. The
carrier must be \emph{regenerated} by removing the modulation.

\subsection{Data-aided and decision-directed detectors}

If the symbols are known, the modulation is removed by multiplying by their conjugates, and the
maximum-likelihood phase estimate~\eqref{eq:ch10:mlphase} over a block is
$\hat\theta=\arg\sum a_k^*z_k$. Its per-symbol version, $e_k=\im\{z_ka_k^*\}$ (or $\arg(z_ka_k^*)$),
is the \term{data-aided phase detector}, with S-curve $\sin\phi$ (or $\phi$ itself). It is the ideal
detector during a preamble or on pilot symbols.

Replacing $a_k$ by the decision $\hat a_k$ gives the \term{decision-directed} detector
$e_k=\im\{z_k\hat a_k^*\}$, which works for any constellation, PSK or QAM, and is what
\texttt{commlib.pll\_dd} implements. When the phase error is small and the SNR high, decisions are
correct and the detector is as good as a data-aided one. But the decision regions of an $M$-PSK
constellation are symmetric under rotation by $2\pi/M$, so the detector's S-curve repeats with period
$2\pi/M$: at a phase error of $\pi/4$ a QPSK decision flips to the neighbouring point and the error
signal changes sign. Figure~\ref{fig:ch10_pd_scurves}(b) shows the result. At high SNR the S-curve is a
sawtooth with period $\pi/2$; as the SNR falls, decision errors near the boundaries flatten it,
reducing the slope (and hence the loop bandwidth) and adding noise.

\mdcfig{ch10_pd_scurves}{Phase-detector S-curves (mean output against phase error), measured by
simulation. (a) The data-aided detector has period $2\pi$; the BPSK Costas detector $2IQ=\sin2\phi$ has
period $\pi$. (b) QPSK decision-directed and hard-limited Costas detectors have period $\pi/2$; noise
flattens the sawtooth, and at 3\,dB the decision-directed detector's slope has almost vanished.}

For QAM the situation is worse: the inner points of 16-QAM or 64-QAM lie closer to their neighbours
in angle than the corners do, so decisions become unreliable at smaller phase errors. A 64-QAM
decision-directed loop is reliable only once the phase error is below a few degrees. Practical QAM
receivers therefore acquire phase with a pilot-aided or \emph{reduced-constellation} detector that
uses only the corner points (whose phases are those of QPSK) and switch to full decision direction
once locked.

\subsection{The Costas loop}

John Costas's loop (1956) is the classic carrier recovery circuit for suppressed-carrier BPSK
(Figure~\ref{fig:ch10_costas}). The received signal is mixed with in-phase and quadrature outputs of
the VCO, each arm is low-pass filtered (in a digital receiver, matched filtered), and the product
of the arms drives the loop. With a BPSK symbol $d=\pm1$ and phase error $\phi$, the arms carry
$I=d\cos\phi$ and $Q=d\sin\phi$, and the product is
\begin{equation}
e=I\cdot Q=d^2\cos\phi\sin\phi=\tfrac12\sin2\phi,
\end{equation}
independent of the data because $d^2=1$. The loop locks at $\phi=0$ or $\phi=\pi$: a two-fold
ambiguity, the price of blindness. For QPSK the equivalent loop combines the arms as
$e=\operatorname{sgn}(I)\,Q-\operatorname{sgn}(Q)\,I$, which is the hard-limited Costas detector in
\texttt{commlib.costas\_qpsk} and in GNU Radio's \texttt{Costas Loop} block; it has a four-fold
ambiguity and, as Figure~\ref{fig:ch10_pd_scurves}(b) shows, is essentially a decision-directed
detector. Costas loops demodulate a great deal of military and space telemetry to this day.

\begin{figure}[htbp]\centering
\begin{tikzpicture}[
  blk/.style={draw=navy,fill=navy!6,thick,minimum height=0.75cm,minimum width=1.3cm,align=center,font=\sffamily\scriptsize},
  sblk/.style={draw=navy,fill=navy!6,thick,minimum height=0.5cm,minimum width=1.0cm,align=center,font=\sffamily\scriptsize},
  mul/.style={draw=navy,thick,circle,inner sep=0.5pt,font=\small},
  arr/.style={-{Stealth[length=2mm]},thick,navy}]
\node (in) at (0,0) {\small $r(t)$};
\node[mul] (mi) at (2.2,1.5) {$\times$};
\node[mul] (mq) at (2.2,-1.5) {$\times$};
\node[sblk] (vco) at (2.2,0.5) {VCO};
\node[sblk] (ph) at (2.2,-0.5) {$-90^\circ$};
\node[blk] (lpi) at (4.2,1.5) {LPF /\\matched};
\node[blk] (lpq) at (4.2,-1.5) {LPF /\\matched};
\node[mul] (pr) at (6.6,0.5) {$\times$};
\node[blk] (lf) at (4.6,0.5) {loop\\filter};
\node[font=\scriptsize,anchor=west] (dout) at (6.0,2.2) {data out};
\draw[thick,navy] (in) -- (0.9,0);
\draw[arr] (0.9,0) |- (mi.west); \draw[arr] (0.9,0) |- (mq.west);
\draw[arr] (mi)--(lpi); \draw[arr] (mq)--(lpq);
\draw[arr] (lpi.east) -| node[pos=0.25,above,font=\scriptsize]{$I$} (pr.north);
\draw[arr] (lpq.east) -| node[pos=0.25,below,font=\scriptsize]{$Q$} (pr.south);
\draw[arr] (pr)--node[above,font=\scriptsize]{$IQ$}(lf);
\draw[arr] (lf)--(vco.east);
\draw[arr] (vco.north)--node[right,font=\scriptsize]{$\cos$}(mi.south);
\draw[arr] (vco.south)--(ph.north);
\draw[arr] (ph.south)--node[right,font=\scriptsize]{$-\sin$}(mq.north);
\draw[arr] (5.4,1.5) |- (dout.west);
\end{tikzpicture}
\caption{The Costas loop for BPSK. The product of the in-phase and quadrature arms, $\tfrac12\sin2\phi$,
does not depend on the data, so the loop can lock to a suppressed carrier; the in-phase arm is the
demodulated data.}
\label{fig:ch10_costas}
\end{figure}

\subsection{Squaring loss and \texorpdfstring{$M$}{M}th-power loops}

An equivalent way to see the Costas loop is that it squares the signal. Squaring BPSK,
$(d\,e^{j\theta})^2=e^{j2\theta}$, removes the data and leaves a pure carrier at twice the phase; a PLL
tracking it and dividing by two recovers $\theta$ modulo $\pi$. For $M$-PSK the \term{$M$th-power
loop} raises the signal to the $M$th power. The nonlinearity is not free. Squaring the noisy signal
multiplies noise by noise as well as signal by noise, and the extra noise-by-noise term costs
performance at low SNR. The penalty is the \term{squaring loss} $S_L$: the phase variance of the
Costas or squaring loop is the data-aided value divided by $S_L$, and for BPSK with matched-filter arms
\begin{equation}
S_L=\frac{1}{1+1/(2\esno)}.
\end{equation}
At $\esno=10$\,dB the loss is only 0.2\,dB, but at 0\,dB it is 1.8\,dB, and for QPSK and higher-order
PSK the losses are larger. This is the quantitative reason that blind carrier recovery struggles
at the low SNRs of heavily coded systems.

\subsection{Feedforward phase estimation}

For bursts, a loop's acquisition transient wastes symbols. The feedforward alternative estimates
the phase over a block. The \term{Viterbi--Viterbi} estimator (1983) generalises the $M$th-power
idea to a block of $N$ symbols:
\begin{equation}
\hat\theta=\frac1M\arg\sum_{k=0}^{N-1}F(|z_k|)\,e^{jM\arg z_k},
\end{equation}
where the nonlinearity $F(\cdot)$ can be chosen to optimise performance at a given SNR ($F(x)=x^M$
recovers the plain $M$th power; $F(x)=x^2$ or $F(x)=1$ is often better for $M=4$ at low SNR). For
QPSK with points at odd multiples of $\pi/4$, $z^4=-|z|^4e^{j4\theta}$, so the estimate is
$\tfrac14\arg(-\sum z_k^4)$. Because each block's estimate lies in $(-\pi/M,\pi/M]$, successive blocks
must be \emph{unwrapped}---whenever the estimate jumps by more than $\pi/M$, add or subtract $2\pi/M$---to
follow a phase that drifts through the ambiguity boundaries. Feedforward estimators are popular in
coherent optical receivers, where the symbol rates of tens of gigabaud make a feedback loop's
latency intolerable and the phase noise of the lasers demands rapid tracking; the blind phase search
algorithm used for 16- and 64-QAM there is a brute-force feedforward ML estimator that tests dozens of
candidate phases in parallel.

\subsection{Phase ambiguity}

Every blind or decision-directed carrier recovery scheme for a symmetric constellation has an
$M$-fold (for square QAM, four-fold) \term{phase ambiguity}: rotating the whole constellation by a
symmetry angle produces an equally valid lock point. Three remedies are in use.
\begin{description}
\item[Differential encoding.] Encode information in the \emph{change} of phase from symbol to
symbol, so a constant rotation cancels. The cost is error multiplication: one wrong phase decision
corrupts two differences, roughly doubling the bit error rate at high SNR (about 0.5\,dB or less in
required SNR, depending on the scheme). DQPSK in TETRA and the differential encoding in many satellite
modems use this approach.
\item[Unique words and preambles.] A known sequence tells the receiver which of the $M$ rotations it
has locked to: correlate against all $M$ rotated versions (or simply take the phase of the complex
correlation peak, as in \texttt{commlib.frame\_sync}) and rotate accordingly.
\item[Rotationally invariant codes.] Trellis-coded modulation and some modern coded schemes are
designed so that rotated codewords are also codewords of a related code, allowing the decoder to
work through the ambiguity with differential encoding of only a subset of bits.
\end{description}
Cycle slips re-open the ambiguity at a random time, which is why continuous-mode systems insert unique
words or pilots periodically rather than only once at the start.

\subsection{Pilot symbols}

The cleanest solution to low SNR, ambiguity and cycle slips alike is to transmit known pilot
symbols in the data stream. In \term{pilot-symbol-assisted modulation} (Cavers, 1991) the transmitter inserts one known
symbol every $P$ data symbols; the receiver estimates the phase (and, in fading channels, the complex
gain) at each pilot and interpolates between them. The overhead is $1/P$, typically 2--5\%, and the
interpolation filter plays the role of the loop filter. The pilots must be dense enough to sample the
fastest phase variation (from residual CFO, phase noise or Doppler) above the Nyquist rate. DVB-S2
inserts optional 36-symbol pilot blocks every 1440 symbols; 5G NR adds phase-tracking reference
signals (PT-RS) at millimetre-wave frequencies, where oscillator phase noise is severe
(Chapter~\ref{ch:ofdm}); and every OFDM system from 802.11a onward reserves a few pilot subcarriers in
every symbol for exactly this purpose.

\begin{history}[Costas and the suppressed carrier, 1956]
John P. Costas of General Electric published ``Synchronous communications'' in the
\emph{Proceedings of the IRE} in December 1956. His argument was economic as much as technical: an
AM transmitter spends at least two-thirds of its power on the carrier, which carries no information,
and a double-sideband suppressed-carrier (DSB-SC) system could put all that power into the
sidebands---if only the receiver could regenerate the carrier. The loop that now bears his name did
exactly that, and Costas showed that synchronous DSB-SC could outperform SSB in some conditions
while being simpler to receive. The loop did not displace SSB for voice radio, but it found its true
home a few years later, when digital phase modulation arrived: every BPSK and QPSK receiver is, in
one form or another, a Costas loop.
\end{history}

\section{Frequency offset estimation}
\label{sec:ch10:freq}

Frequency is the synchronization parameter that most often defeats a receiver. A phase error is
bounded; a frequency error makes the phase grow without limit, and once the accumulated rotation
over the matched filter's span or a correlation window becomes large, every other estimator in the
chain fails. The design problem is a trade-off between \emph{range} and \emph{accuracy}, and it is
almost always solved in two steps: coarse then fine.

\subsection{The range--accuracy trade-off}

Most frequency estimators measure the phase advance between samples separated by a lag $D$:
\begin{equation}
\hat\nu=\frac{1}{2\pi D}\arg\sum_k z_k\,z_{k-D}^*,
\label{eq:ch10:delaymult}
\end{equation}
where $z_k$ is the signal after the modulation has been removed (by known symbols or a nonlinearity)
and $\nu$ is in cycles per sample. The phase advance $2\pi\nu D$ can only be measured modulo $2\pi$, so
the unambiguous range is $|\nu|<1/(2D)$. But the accuracy improves with $D$, because the same phase
noise is divided by a larger lever arm. A short lag gives a wide range and poor accuracy; a long lag
the reverse. The 802.11 preamble (Section~\ref{sec:ch10:frame}) is built around exactly this
trade-off: a short training field with lag 16 samples for coarse estimation, then a long training
field with lag 64 samples for fine estimation.

\subsection{Non-data-aided: the \texorpdfstring{$M$}{M}th-power spectrum}

Without known symbols, the modulation of $M$-PSK is removed by the $M$th power: raising
$a_ke^{j(2\pi\nu k+\theta)}$ to the power $M$ leaves a pure tone at $M\nu$ plus noise. An FFT of $z_k^M$ therefore shows a spectral line at $M\nu$
(Figure~\ref{fig:ch10_mpower}(a)), and the peak location divided by $M$ is the estimate. This is the
function \texttt{cfo\_power\_estimate} in \texttt{commlib.sync}, and it is a popular coarse estimator for satellite burst
modems. Its unambiguous range is $|\nu|<1/(2M)$ cycles/symbol---for QPSK, one eighth of the symbol
rate---beyond which the estimate aliases (Figure~\ref{fig:ch10_mpower}(b)). The $M$th power also
multiplies the noise, so the estimator's threshold is high: at 3\,dB the $N=256$ estimate in the
figure already produces outliers.

\mdcfig{ch10_mpower}{The fourth-power frequency estimator for QPSK. (a) The spectrum of the received
signal shows no line, but its fourth power has a strong line at $4\nu$. (b) Estimates against true
offset: the range is limited to $\pm1/8$ cycle per symbol, and at low SNR outliers appear even inside
it.}

\subsection{Data-aided estimators}

When the symbols are known (a preamble or pilots), multiplying by $a_k^*$ leaves
$z_k=e^{j(2\pi\nu k+\theta)}+\text{noise}$: a single complex tone in white noise, the most studied
estimation problem in signal processing. Its maximum-likelihood solution, shown by Rife and Boorstyn
in 1974, is the location of the peak of the periodogram $|\sum_kz_ke^{-j2\pi\nu k}|^2$, found by a
zero-padded FFT followed by interpolation or a fine search. It meets the CRB above threshold and has a
range of $\pm1/2$. Its cost is an FFT and a search; several cheaper estimators based on
the autocorrelation $R(m)=\frac{1}{N-m}\sum_{k=m}^{N-1}z_kz_{k-m}^*$ come close:
\begin{description}
\item[Kay (1989)] averages the phase differences of adjacent samples with a parabolic window,
\[
\hat\nu=\frac{1}{2\pi}\sum_{k=1}^{N-1}w_k\arg(z_kz_{k-1}^*),\qquad
w_k=\frac{3N/2}{N^2-1}\Bigl[1-\Bigl(\frac{2k-N}{N}\Bigr)^2\Bigr].
\]
It has full range $\pm1/2$ and meets the
CRB at high SNR, but its threshold is high, because each phase difference is computed from only two
noisy samples.
\item[Fitz (1994)] uses the phases of $R(m)$ for $m=1,\dots,L$:
\[
\hat\nu=\frac{1}{\pi L(L+1)}\sum_{m=1}^{L}\arg R(m).
\]
Averaging before taking the phase lowers the
threshold dramatically, at the cost of range $|\nu|<1/(2L)$.
\item[Luise--Reggiannini (1995)] sums the correlations first and takes one phase:
\[
\hat\nu=\frac{1}{\pi(L+1)}\arg\sum_{m=1}^{L}R(m),
\]
with range $|\nu|<1/(L+1)$. With $L=N/2$ it is
within a fraction of a decibel of the CRB and has a low threshold; it is the default choice for
burst-mode preambles.
\end{description}
Figure~\ref{fig:ch10_freq_est} compares all four on a 64-symbol preamble. All approach the CRB at high
SNR; the differences are in the threshold, which spans more than 12\,dB, and in the range. The
autocorrelation estimators, notice, are all instances of the delay-and-multiply
estimator~\eqref{eq:ch10:delaymult} with different lags averaged in different ways.

\mdcfig[0.7]{ch10_freq_est}{Mean-square error of data-aided frequency estimators on a 64-symbol
preamble ($\nu$ uniform in $\pm0.01$). All reach the Cram\'er--Rao bound at high SNR, but Kay's
estimator breaks down below about 10\,dB, the periodogram below $-2$\,dB, while Fitz and
Luise--Reggiannini (with $L=N/2$, and a correspondingly reduced range) remain on the bound
throughout.}

\begin{table}[!htb]\centering\small
\caption{Data-aided frequency estimators on $N$ symbols. Ranges in cycles/symbol.}
\label{tab:ch10:freqest}
\begin{tabular}{llll}
\toprule
Estimator & Range & Threshold & Complexity \\
\midrule
Periodogram (Rife--Boorstyn, ML) & $\pm1/2$ & low & FFT + search \\
Kay & $\pm1/2$ & high & $N$ phase computations \\
Fitz, lags $1\dots L$ & $\pm1/(2L)$ & low & $L$ correlations, $L$ phases \\
Luise--Reggiannini, lags $1\dots L$ & $\pm1/(L+1)$ & low & $L$ correlations, one phase \\
Delay-and-multiply, lag $D$ & $\pm1/(2D)$ & moderate & one correlation \\
\bottomrule
\end{tabular}
\end{table}

\subsection{Frequency-locked loops}

A \term{frequency-locked loop} (FLL) is a feedback loop whose detector measures frequency error
rather than phase error. The classic \emph{cross-product} discriminator computes
$e_k=\im\{z_kz_{k-1}^*\}$, proportional to $\sin(2\pi\Delta f T)$, which is insensitive to the absolute
phase. An FLL has no phase ambiguity, no hang-up and a wide pull-in range; it is less accurate than a
PLL and leaves a random phase, so it is used for acquisition and then hands over to a PLL, or runs
alongside it as an \emph{FLL-assisted PLL}, the standard architecture in GNSS receivers, where
Doppler dynamics can be severe. For PSK signals the modulation can be removed before the
discriminator (a $4$th-power or decision-directed FLL), or the \emph{band-edge} FLL of fred harris
can be used: it compares the energy in two filters matched to the upper and lower band edges of the
RRC spectrum, which are equal only when the signal is centred. GNU Radio's \texttt{FLL Band-Edge}
block implements it, and it is the usual first stage of a GNU Radio PSK receiver.

\begin{inpractice}[Doppler in low Earth orbit]
A satellite in a 550\,km orbit moves at about 7.6\,km/s. At a Ku-band downlink of 12\,GHz the maximum
Doppler shift approaches $v/c\cdot f\approx300$\,kHz, and it sweeps through zero in the few minutes of a pass, with
a peak rate of a few kilohertz per second overhead. A terminal must therefore acquire over several
hundred kilohertz and track a large Doppler rate, which a type-2 loop can follow only with a phase lag
$\dot\omega/\omega_n^2$ (Table~\ref{tab:ch10:steady}). Real terminals predict the Doppler from the
satellite ephemeris and their own position and pre-compensate, leaving only a small residual for the
loops. 5G non-terrestrial networks (3GPP Release 17) take the same approach: the phone uses
broadcast ephemeris and its GNSS position to pre-compensate uplink frequency and timing before
transmitting (Chapter~\ref{ch:satellite}).
\end{inpractice}

\subsection{Frequency synchronization in OFDM}

OFDM (Chapter~\ref{ch:ofdm}) is far more sensitive to frequency offset than single-carrier
modulation, because its subcarriers are closely spaced and orthogonal only when the receiver's
FFT bins line up exactly with them. A CFO of a fraction $\epsilon$ of the subcarrier spacing rotates
every subcarrier by a common phase and also causes \emph{inter-carrier interference} (ICI) whose power
is roughly $(\pi\epsilon)^2/3$ of the signal power for small $\epsilon$. Pollet, Van Bladel and
Moeneclaey showed that the resulting SNR degradation is approximately
\begin{equation}
D\approx\frac{10}{3\ln10}\,(\pi\epsilon)^2\,\esno\ \text{dB},
\end{equation}
so at $\esno=20$\,dB an offset of only 1\% of the subcarrier spacing costs about 0.15\,dB, and 2\%
about 0.6\,dB. OFDM
receivers split the offset into a \emph{fractional} part (less than half a subcarrier), estimated from
the phase of a correlation between repeated time-domain segments, and an \emph{integer} part (a whole
number of subcarriers), estimated by finding the shift of a known frequency-domain sequence.
Moose (1994) derived the ML estimator for two identical OFDM symbols; Schmidl and Cox (1997)
designed a preamble whose first symbol has two identical halves, so that the same correlation that
finds the frame gives the fractional offset (Section~\ref{sec:ch10:frame}). The CP itself is a
repeated segment, which is exploited by the blind estimator of van de Beek, Sandell and B\"orjesson
(1997).

\begin{worked}[Frequency acquisition in Wi-Fi]
IEEE 802.11 allows each device a carrier frequency tolerance of $\pm20$\,ppm, so two devices at
5.8\,GHz can differ by up to 40\,ppm, or 232\,kHz. The short training field repeats every
$0.8\,\mu$s (16 samples at 20\,MS/s); by~\eqref{eq:ch10:delaymult} its delay-and-correlate estimate is
unambiguous for $|\Delta f|<1/(2\times0.8\,\mu\text{s})=625$\,kHz, comfortably covering the worst case. The
long training field repeats every $3.2\,\mu$s, with range $\pm156$\,kHz: not enough on its own, but
after coarse correction the residual is a few kilohertz, and the long lag gives four times the
accuracy per unit of noise. A receiver that used only the long training field would mistake a
200\,kHz offset for $-112.5$\,kHz and fail to decode, which is precisely why the preamble has two
parts.
\end{worked}

\section{Symbol timing recovery}
\label{sec:ch10:timing}

The matched filter of Chapter~\ref{ch:baseband} maximises the SNR only at one instant per symbol,
and a Nyquist pulse is free of ISI only at those instants. Sampling a quarter of a symbol late costs
several decibels (Figure~\ref{fig:ch10_timing_ber}). Symbol timing recovery---also called clock
recovery or symbol synchronization---finds those instants and follows them as the transmitter's clock
and the receiver's drift apart.

\subsection{Two architectures}

There are two ways to make the samples land in the right place. In a \emph{synchronous} receiver the
timing loop steers the ADC's sampling clock itself, through a voltage-controlled crystal oscillator,
so that the samples arrive at the right instants. This was the norm in early digital modems and
survives in some SerDes and in CDR circuits (Section~\ref{sec:ch10:cdr}). In an
\emph{asynchronous} (or \emph{free-running}) receiver the ADC samples at a fixed rate unrelated to the
symbol rate, typically two to eight samples per symbol, and the timing correction is done entirely in
DSP by \term{interpolation}: computing the value the signal \emph{would have had} at the desired
instant from the samples around it. Every software-defined radio works this way, because the sample
clock is shared by all the signals in the band and cannot be steered for any one of them.

Gardner's framework (Figure~\ref{fig:ch10_timing_loop}) organises the asynchronous timing loop into
four blocks: an \emph{interpolator} that computes $y(kT+\hat\tau)$ from samples $x[n]$; a
\term{timing error detector} (TED) that produces an error signal from the interpolants; a loop filter;
and an \emph{interpolation controller}---an NCO that decides which input sample $m_k$ and which
fractional delay $\mu_k\in[0,1)$ each output symbol needs.

\begin{figure}[htbp]\centering
\begin{tikzpicture}[
  blk/.style={draw=navy,fill=navy!6,thick,minimum height=0.9cm,minimum width=1.55cm,align=center,font=\sffamily\scriptsize},
  arr/.style={-{Stealth[length=2mm]},thick,navy}]
\node (in) at (0,0) {\scriptsize $x[n]$ at $T_s$};
\node[blk] (mf) at (2.0,0) {matched\\filter};
\node[blk] (ip) at (4.4,0) {interpolator\\(Farrow)};
\node[blk] (ted) at (7.0,0) {timing error\\detector};
\node[blk] (lf) at (7.0,-1.6) {loop\\filter};
\node[blk] (nco) at (4.4,-1.6) {interpolation\\control (NCO)};
\node[font=\scriptsize,anchor=west] (out) at (8.3,0.9) {symbols $y(kT+\hat\tau)$};
\draw[arr] (in)--(mf); \draw[arr] (mf)--(ip); \draw[arr] (ip)--(ted);
\draw[arr] (ted)--node[right,font=\scriptsize]{$e_k$}(lf); \draw[arr] (lf)--node[above,font=\scriptsize]{$v_k$}(nco);
\draw[arr] (nco)--node[right,font=\scriptsize,align=left]{$m_k$, $\mu_k$\\strobe}(ip);
\draw[arr] (5.7,0) |- (out.west);
\end{tikzpicture}
\caption{Gardner's asynchronous timing loop. The ADC runs at a fixed rate; the NCO decides when a
symbol strobe occurs and with what fractional delay $\mu_k$ the interpolator must compute it.}
\label{fig:ch10_timing_loop}
\end{figure}

The controller is a modulo-1 counter decremented by $W\approx T_s/T$ (the nominal ratio of sample to
symbol period, corrected by the loop filter output $v_k$) at each input sample. When it underflows, a
symbol strobe is due between the current and the next sample, and the fractional position, computed
from the counter's residue, is $\mu_k$. The loop is a PLL in which the ``phase'' is the timing
offset, and all the loop design of Section~\ref{sec:ch10:pll} carries over: a type-2 loop tracks a
constant clock offset (a timing ramp) with zero steady-state error, and its noise bandwidth $B_LT$
trades acquisition speed against timing jitter.

\subsection{Timing error detectors}

Maximising the likelihood~\eqref{eq:ch10:llf} over $\tau$ leads to the \emph{ML timing error
detector}, which correlates the matched-filter output with its derivative:
$e_k=\re\{\hat a_k^*\,\dot y(kT+\hat\tau)\}$, or blindly $\re\{y^*(kT+\hat\tau)\,\dot y(kT+\hat\tau)\}$.
The derivative is zero at the peak of each pulse, so the detector drives the sampling instant to the
peak. Practical TEDs approximate the derivative with differences between samples.

\paragraph{Early--late gate.} Take samples a little before and after the current estimate, at
$\pm\Delta$ (typically $\Delta=T/4$ or $T/2$), and compare their energies:
$e_k=|y(kT+\hat\tau-\Delta)|^2-|y(kT+\hat\tau+\Delta)|^2$ (the sign chosen so the loop moves toward the
larger). The early--late detector is the oldest and most intuitive; it needs at least two or three
samples per symbol and is used in spread-spectrum code tracking (Chapter~\ref{ch:spreadspectrum}),
where the ``symbol'' is a chip of a pseudo-noise code and the early--late delay-locked loop tracks the
correlation peak.

\paragraph{Mueller and M\"uller.} Mueller and M\"uller (1976) derived a family of detectors that
need only \emph{one sample per symbol}, which mattered greatly when every ADC sample was expensive.
The most widely used member is
\begin{equation}
e_k=\re\{\hat a_{k-1}^*\,y_k-\hat a_k^*\,y_{k-1}\},
\end{equation}
where $y_k$ are the symbol-rate samples and $\hat a_k$ the decisions. Its S-curve
(Figure~\ref{fig:ch10_ted_scurves}(a)) is nearly linear over $\pm T/2$, but it is decision-directed and
therefore needs the carrier phase to be approximately correct, and its performance depends on the
pulse shape and on ISI. It was the standard in voiceband modems and survives in SerDes receivers
(where \emph{baud-rate} CDR saves power at 50 to 100\,GBd) and in GNU Radio's \texttt{Symbol Sync}
block.

\paragraph{Gardner.} Gardner's detector (1986) uses two samples per symbol---the symbol-centre
samples $y_k$, $y_{k-1}$ and the midpoint sample $y_{k-1/2}$ between them:
\begin{equation}
e_k=\re\bigl\{y_{k-1/2}^*\,\bigl(y_k-y_{k-1}\bigr)\bigr\}.
\label{eq:ch10:gardner}
\end{equation}
The intuition is transparent for BPSK. If a transition occurs between symbols $k-1$ and $k$, the
midpoint sample should be zero when timing is correct. If sampling is late, the midpoint has already
moved toward the new symbol, and its sign times the direction of the transition $y_k-y_{k-1}$ gives a
positive error. When no transition occurs, $y_k-y_{k-1}\approx0$ and the detector is silent. Gardner's
detector does not need decisions and is \emph{independent of the carrier phase}: rotating all samples
by $e^{j\phi}$ leaves~\eqref{eq:ch10:gardner} unchanged. That is why timing recovery can run before
carrier recovery in the chain of Figure~\ref{fig:ch10_chain}. It is \texttt{commlib.gardner\_sync} and
the default TED in most software radio receivers.

Its weakness appears in Figure~\ref{fig:ch10_ted_scurves}(b). The detector relies on the zero crossings
between symbols, and with a small roll-off the pulse tails are long and the zero crossings are
poorly defined. The S-curve slope---and hence the loop's signal-to-self-noise ratio---falls roughly in
proportion to $\beta$. For $\beta=0.1$ the Gardner detector is nearly useless, and systems with very
sharp spectra (DVB-S2X uses roll-offs down to 0.05) use data-aided or ML-based timing detectors
instead.

\mdcfig{ch10_ted_scurves}{Timing error detector S-curves. (a) Early--late, Gardner and
Mueller--M\"uller detectors for RRC $\beta=0.35$, normalised to their peak (early--late and Gardner
coincide in shape). (b) The Gardner detector's slope at the lock point shrinks with the roll-off
factor, because the zero crossings it exploits become ill defined.}

\begin{table}[!htb]\centering\small
\caption{Properties of common timing error detectors.}
\label{tab:ch10:ted}
\begin{tabular}{lcccl}
\toprule
Detector & Samples/symbol & Decisions & Carrier phase & Typical use \\
\midrule
Early--late gate & $\ge2$ & no & no & spread spectrum, GNSS \\
Mueller--M\"uller & 1 & yes & yes & SerDes, voiceband, cable \\
Gardner & 2 & no & no & SDR, satellite modems \\
Zero-crossing & 2 & yes & yes & low-cost modems \\
ML (derivative filter) & $\ge2$ & optional & optional & polyphase clock sync \\
\bottomrule
\end{tabular}
\end{table}

\subsection{Interpolation}

The interpolator must compute $y(t)$ at an arbitrary fractional delay $\mu$ between samples. The ideal
interpolator is a sinc, which is infinitely long, so practical interpolators are short FIR filters
whose coefficients depend on $\mu$. Recomputing the coefficients for every symbol would be expensive,
but if they are polynomials in $\mu$ the filter can be written as a \term{Farrow structure}: a small
bank of fixed FIR filters whose outputs are combined by Horner's rule,
\begin{equation}
y(\mu)=\bigl((c_3\mu+c_2)\mu+c_1\bigr)\mu+c_0,
\end{equation}
where each $c_i$ is a fixed four-tap filter applied to the samples $x[m_k-1],\dots,x[m_k+2]$. The cubic
Lagrange interpolator in \texttt{commlib.interp\_cubic} is exactly this, and Erup, Gardner and Harris
(1993) showed that a cheaper \emph{piecewise-parabolic} interpolator with one free parameter performs
almost as well, which made it a favourite in FPGA modems.

Figure~\ref{fig:ch10_interp} shows why oversampling matters. A four-tap interpolator is nearly
perfect at low frequencies but droops and delays unevenly near half the sample rate, and the error is
largest at $\mu=0.5$. With four samples per symbol, the signal occupies only the lower quarter of the
band and a cubic interpolator introduces negligible error; at two samples per symbol the signal
reaches $0.34$ cycles/sample (for $\beta=0.35$) and the interpolation error becomes a measurable
source of self-noise. Linear interpolation (dotted) is adequate only at high oversampling.

\mdcfig{ch10_interp}{Interpolation for timing recovery. (a) The interpolator computes strobes
(diamonds) between the available samples at fractional positions $\mu_k$. (b) Magnitude response of
the cubic Lagrange (Farrow) interpolator for three values of $\mu$, with linear interpolation dotted.
(c) Its group delay, which should equal $\mu$ at all frequencies.}

\subsection{Polyphase filter-bank clock recovery}

An elegant alternative, due to harris and Rice (2001), merges the matched filter and the interpolator.
The matched filter is designed at a high oversampling rate---say 32 times the input rate---and split
into 32 polyphase subfilters (Chapter~\ref{ch:dsp}), each of which is a matched filter with a different
fractional delay. The timing loop simply selects which subfilter to use for each symbol. A second
bank, built from the derivative of the matched filter, supplies $\dot y$ for the ML timing
detector $\re\{y^*\dot y\}$. No separate interpolator is needed, and the fractional-delay resolution
is $1/32$ of a sample. This is the algorithm in GNU Radio's \texttt{Polyphase Clock Sync} block,
the timing recovery in most GNU Radio PSK receivers.

\subsection{Feedforward timing: Oerder and Meyr}

For bursts, a feedforward timing estimator avoids the loop transient. Oerder and Meyr (1988) observed
that the squared magnitude of a linearly modulated signal, $|x(t)|^2$, is cyclostationary with period
$T$: its mean varies periodically, peaking at the symbol instants. Its Fourier series at $1/T$
therefore has a phase that reveals the timing. With $N_s$ samples per symbol (at least 3, typically 4)
and $L_0$ symbols,
\begin{equation}
\hat\tau=-\frac{T}{2\pi}\arg\sum_{n=0}^{L_0N_s-1}|x[n]|^2\,e^{-j2\pi n/N_s}.
\end{equation}
The estimator is blind, independent of carrier phase and frequency offset, and costs a handful of
operations per sample. Figure~\ref{fig:ch10_timing_loops}(b) shows its accuracy against the MCRB. At
low SNR it is within a few decibels of the bound, but at high SNR it reaches a floor set by
\emph{self-noise}: the random data make $|x(t)|^2$ fluctuate even without noise, and the fluctuation is
worse for small roll-off, where the spectral line at $1/T$ is weak.

\mdcfig{ch10_timing_loops}{(a) The Gardner loop (\texttt{commlib.gardner\_sync}) acquiring a
1.7-sample delay, and tracking a 200\,ppm clock offset as a steady ramp of the strobe phase (unwrapped).
(b) Oerder--Meyr feedforward timing estimation over 100 symbols: mean-square error against the MCRB
(dashed). Self-noise from the random data sets an error floor that is worst for small roll-off.}

\begin{worked}[How fast does the clock drift?]
Two USRP B200 radios, each with a $\pm2$\,ppm TCXO, may differ by 4\,ppm. At 1\,Msym/s the timing error
grows by $4\times10^{-6}\times10^6=4$ symbol periods per second, or one symbol every 250\,ms. A 1\,ms burst
drifts by only $0.004T$ and a one-shot timing estimate from the preamble suffices. A 100\,ms capture
drifts by $0.4T$---enough to destroy the eye---and needs a tracking loop. For the loop to track a clock
offset $\epsilon_s$ with zero steady-state error it must be type 2; its integrator converges to
$\epsilon_s$, which is also an excellent estimate of the carrier frequency offset if both oscillators
are derived from the same reference, as they are in the B200. Many receivers exploit this to aid the
carrier loop with the timing loop's frequency estimate, or vice versa.
\end{worked}

\subsection{Timing jitter and its cost}

The timing loop's residual error---a static offset from detector bias plus random jitter from
noise and self-noise---degrades detection. Figure~\ref{fig:ch10_timing_ber} quantifies it for a raised
cosine channel: with $\beta=0.35$, a static offset of $0.05T$ costs QPSK about 0.4\,dB at a BER of
$10^{-5}$ and $0.1T$ costs about 1.4\,dB, while 16-QAM, with its smaller decision distances, loses about
1.7\,dB at $0.05T$. Smaller roll-offs are more sensitive, because their pulses' tails are larger and the
ISI at a given offset greater. A useful design target for high-order QAM is an RMS timing jitter of
1--2\% of the symbol period, which, from the loop-noise relation, dictates a timing loop bandwidth
$B_LT$ of the order of $10^{-3}$ at typical SNRs.

\mdcfig{ch10_timing_ber}{The cost of timing error with raised-cosine pulses, computed by averaging
exact conditional error probabilities over random ISI patterns. (a) QPSK bit error rate with static
offsets. (b) The $\ebno$ penalty at BER $10^{-5}$: 16-QAM (dashed) is several times more sensitive than
QPSK, and small roll-offs are more sensitive than large ones.}

\begin{inpractice}[Timing recovery in GNU Radio]
GNU Radio provides three generations of symbol synchronizers, and knowing which to use saves hours.
The \texttt{Clock Recovery MM} block (Mueller--M\"uller with a simple interpolator) is the oldest,
and the two newer blocks supersede it for most purposes. The \texttt{Polyphase Clock Sync} block implements harris and Rice's filter-bank
method with an ML detector; it combines matched filtering and timing and is the usual choice for PSK.
The \texttt{Symbol Sync} block is the most general: it offers Gardner, Mueller--M\"uller, early--late,
zero-crossing and ML detectors, several interpolators, and a loop parameterised directly by $B_LT$ and
$\zeta$, exactly as in this chapter. A typical QPSK receiver chain is \texttt{FLL Band-Edge}
$\rightarrow$ \texttt{Symbol Sync} (or \texttt{Polyphase Clock Sync}) $\rightarrow$ \texttt{Costas Loop}
$\rightarrow$ decision. The companion lab reproduces this chain with \texttt{commlib}.
\end{inpractice}

\section{Frame synchronization}
\label{sec:ch10:frame}

Carrier and timing recovery deliver a stream of correct symbols; frame synchronization tells the
receiver what they mean. It must find the beginning of each packet, frame or codeword, and in burst
systems it must also detect that a packet is present at all. It is the last synchronizer in the chain
and, because it usually rests on a known sequence, it often supplies the data-aided references that
the earlier loops lack.

\subsection{Correlating with a known sequence}

The transmitter inserts a known sequence $s_0,\dots,s_{N-1}$---a \term{preamble} at the start of a burst,
or a periodically repeated \term{unique word} or \emph{sync word} in a continuous stream. The receiver
correlates:
\begin{equation}
c(\mu)=\sum_{i=0}^{N-1}s_i^*\,y_{i+\mu},
\end{equation}
and declares the frame start at the lag $\mu$ where $|c(\mu)|$ exceeds a threshold or is largest. Using
the magnitude makes the correlation \emph{non-coherent}: it works before the carrier phase is known,
and the phase of $c(\hat\mu)$ is then the data-aided phase estimate that resolves the carrier loop's
ambiguity. \texttt{commlib.frame\_sync} does precisely this, normalising $|c|$ by the energy in the
window so that the threshold does not depend on the signal level.

A residual frequency offset $\nu$ rotates the sequence during the correlation and reduces the peak
by the factor $|\sin(\pi\nu N)/(N\sin\pi\nu)|$; the peak falls by 4\,dB when $\nu N=0.5$. Long sequences
in the presence of large offsets therefore use \emph{segmented} correlation---correlate coherently over
$N/P$ symbols, then add the $P$ segment magnitudes non-coherently---or \emph{differential} correlation,
which correlates $y_{i}y_{i-1}^*$ against $s_is_{i-1}^*$ and is nearly immune to frequency offset at the
cost of some noise enhancement, since each product contains two noisy samples.

\subsection{Sequences with good correlation}

The ideal sync sequence has a sharp autocorrelation peak and low sidelobes, so that no shifted
version of it, and no data, looks like the real thing. Three families dominate
(Figure~\ref{fig:ch10_sequences}).
\begin{description}
\item[Barker sequences] are binary sequences whose \emph{aperiodic} autocorrelation sidelobes are at
most 1 in magnitude. They exist only for lengths 2, 3, 4, 5, 7, 11 and 13, and none longer is known;
Barker described them in 1953. The length-11 Barker code spreads each bit of 1 and 2\,Mb/s 802.11b
Wi-Fi, and the length-13 code is a classic radar pulse compression code.
\item[Maximal-length sequences] (m-sequences), generated by a linear feedback shift register with a
primitive polynomial, have length $2^m-1$ and a two-valued \emph{periodic} autocorrelation: $N$ at zero
lag and $-1$ everywhere else. Their aperiodic sidelobes are larger, of order $\sqrt N$, which matters
for one-shot preambles. Gold and Kasami sets, built from pairs of m-sequences, trade a little
autocorrelation for low cross-correlation between members, which is what GPS and CDMA need
(Chapter~\ref{ch:spreadspectrum}).
\item[Zadoff--Chu sequences], $s_n=e^{-j\pi un(n+1)/N}$ for odd length $N$ and root $u$ coprime to $N$,
are complex with constant amplitude and \emph{zero} periodic autocorrelation sidelobes (a
\emph{CAZAC} sequence). Different roots have constant cross-correlation magnitude $\sqrt N$. Their
constant envelope keeps the power amplifier happy, and their DFT is also a Zadoff--Chu sequence, so
they are equally good in time and frequency. LTE uses them for its primary synchronization signal, its
random-access preambles and its uplink reference signals, and NR for random access and reference
signals.
\end{description}

\mdcfig{ch10_sequences}{Correlation properties of synchronization sequences. (a) Barker-13: aperiodic
sidelobes of at most 1. (b) An $N=63$ m-sequence: periodic autocorrelation $-1$ off-peak (dark), but
larger aperiodic sidelobes (grey). (c) Zadoff--Chu: zero periodic sidelobes, and a flat
cross-correlation of $\sqrt{63}\approx7.9$ between roots.}

\subsection{Threshold design}

In a burst receiver, the frame synchronizer is also the packet detector, and its threshold sets two
error probabilities: \emph{false alarm} (declaring a packet in noise or data) and \emph{miss} (failing
to see a real one). Normalise the correlation as $\Lambda(\mu)=|c(\mu)|^2/(N\sigma^2)$, where $\sigma^2$
is the power of the signal-plus-noise at a lag where the sequence is absent. Off the peak, $c$ is a sum
of many random terms and approximately complex Gaussian, so $\Lambda$ is exponentially distributed and
the per-lag false-alarm probability is
\begin{equation}
P_{FA}=e^{-\Lambda_t}
\end{equation}
for threshold $\Lambda_t$. On the peak, $c$ has a deterministic component $N\sqrt{E_s}$ plus noise, and
$\Lambda$ follows a non-central chi-square distribution; the miss probability is a Marcum $Q$-function
that depends on the \emph{total} preamble energy $NE_s/N_0$. Figure~\ref{fig:ch10_frame_detect} plots
both. The false-alarm curve does not depend on $N$ or SNR at all; the miss curves move right as $N$
grows. Choosing $N$ and $\Lambda_t$ is a matter of finding a threshold that sits between them.

\mdcfig{ch10_frame_detect}{Frame detection with a correlation threshold. (a) The normalised correlation
of a 32-symbol QPSK preamble embedded in random data at $\esno=0$\,dB, with the threshold for a
$10^{-3}$ false-alarm probability over a 600-lag search. Note the data-induced peak near lag 165.
(b) Per-lag false-alarm probability and miss probability against threshold for $N=16$, 32 and 64
(solid 0\,dB, dashed $-3$\,dB); circles are a Monte Carlo check.}

\begin{worked}[Sizing a preamble]
A burst receiver searches a window of 600 lags for a QPSK preamble at $\esno=0$\,dB and must keep
the overall false-alarm probability per window below $10^{-3}$. With a union bound, the per-lag
$P_{FA}$ must be $10^{-3}/600$, so $\Lambda_t=\ln(6\times10^5)=13.3$. From the non-central chi-square
distribution, a 32-symbol preamble then misses 22\% of packets---unacceptable. Doubling to 64 symbols
reduces the miss probability to about $2\times10^{-5}$, a thousandfold improvement for 3\,dB more
preamble energy. If the SNR can drop to $-3$\,dB, even 64 symbols miss 4\% of packets. This steep
dependence on total preamble energy $NE_s/N_0$ is why low-SNR systems (satellite return links, LoRa,
NB-IoT) use very long preambles, and why the preamble is sized at the lowest operating SNR, not the
typical one.
\end{worked}

\subsection{Optimum frame synchronization}

Correlation is optimal when the sequence is surrounded by nothing, but a unique word embedded in a
stream of random data is a different problem: the data around it also contribute to the correlation
and occasionally mimic the sync word. Massey (1972) derived the optimal rule for binary antipodal
signalling in AWGN. For received samples $r_i$, sync bits $s_i\in\{\pm1\}$ and amplitude $\sqrt{E_s}$,
the ML frame position maximises
\begin{equation}
L(\mu)=\sum_{i=0}^{N-1}\Bigl[\frac{2\sqrt{E_s}}{N_0}\,s_i\,r_{i+\mu}-\ln\cosh\Bigl(\frac{2\sqrt{E_s}}{N_0}\,r_{i+\mu}\Bigr)\Bigr],
\end{equation}
which at high SNR becomes $\sum_i\bigl(s_ir_{i+\mu}-|r_{i+\mu}|\bigr)$: the correlation \emph{minus a
correction} for the energy of the samples. The correction penalises lags where the received samples
are large but merely happen to agree with the sync word. Massey's rule gives a useful improvement
over plain correlation, especially for short sync words, and costs almost nothing.

\subsection{Frame synchronization in continuous streams}

A continuous stream such as a T1 line, a SONET link or a satellite broadcast carries a sync word
at a fixed period. The receiver need not detect it perfectly every time, because it knows where the
next one should be. Frame synchronizers are therefore \emph{state machines}: in the \emph{search}
state the receiver looks for the sync word everywhere; after a detection it enters a \emph{verify}
state and checks for the word at the expected position in the next few frames; after enough
confirmations it declares \emph{lock} and thereafter tolerates several consecutive misses (a
\emph{flywheel}) before declaring loss of frame. The numbers of confirmations and tolerated misses
trade acquisition time against robustness. SONET and SDH frame on the bytes A1\,A2 (hexadecimal F6
and 28) at the start of every 125\,$\mu$s frame; T1 frames on a 12-bit pattern spread across a
superframe (Chapter~\ref{ch:pcm}); and every codeword-based system also has the option of using the
decoder itself as a sync detector, since a correctly aligned codeword passes its parity checks and a
misaligned one does not.

\begin{inpractice}[The DVB-S2 physical-layer header]
DVB-S2 (ETSI EN 302 307-1) is a continuous broadcast but is built from self-synchronizing frames. Each
physical-layer frame begins with a 90-symbol PLHEADER modulated in $\pi/2$-BPSK: a 26-symbol
\emph{start-of-frame} (SOF) field, the fixed pattern 18D2E82 in hexadecimal, followed by a 64-symbol
PLS code that carries the modulation, code rate, frame length and pilot flag, protected by a
first-order Reed--Muller code. The payload is divided into slots of 90 symbols, and when pilots are
enabled a block of 36 known symbols is inserted after every 16 slots. A receiver typically runs a
differential correlator for the SOF, which tolerates the large residual frequency offsets of a
satellite downlink, then estimates the frequency from the header and pilots with a data-aided
estimator such as Luise--Reggiannini, tracks the phase from pilot to pilot, and uses the PLS code to
learn how to demodulate the frame that follows. The $\pi/2$-BPSK header can be received at SNRs well
below those needed by the payload, so the receiver can always find the frames, even when it cannot
yet decode them.
\end{inpractice}

\subsection{OFDM preambles: Schmidl--Cox and 802.11}

OFDM receivers usually find the frame and the fractional CFO together with a \emph{delay-and-correlate}
metric that exploits a repeated segment. Schmidl and Cox (1997) proposed a training symbol whose two
halves, of length $L=N/2$, are identical (made by loading only the even subcarriers). Their timing
metric is
\begin{equation}
M(d)=\frac{|P(d)|^2}{R(d)^2},\qquad P(d)=\sum_{m=0}^{L-1}y_{d+m}^*\,y_{d+m+L},\quad R(d)=\sum_{m=0}^{L-1}|y_{d+m+L}|^2.
\end{equation}
$M(d)$ approaches one when the window straddles the two identical halves and is small elsewhere,
regardless of the channel, because a multipath channel affects both halves in the same way. The
phase of $P(\hat d)$ is $\pi\epsilon$ for a CFO of $\epsilon$ subcarrier spacings, so the same computation
gives the fractional frequency offset with range $\pm1$ subcarrier. Because the cyclic prefix is also a
copy, the metric has a \emph{plateau} the length of the CP (Figure~\ref{fig:ch10_ofdm_timing}(a)), which
blurs the timing estimate; Schmidl and Cox take the midpoint of the 90\% points, and later refinements
use a cross-correlation with the known preamble for fine timing.

The legacy IEEE 802.11a/g preamble (Figure~\ref{fig:ch10_ofdm_timing}(b)) is the canonical practical
design. Its \emph{short training field} (STF) is ten repetitions of a 0.8\,$\mu$s pattern, made by
loading only every fourth subcarrier; a delay-and-correlate detector with a 16-sample lag produces a
long plateau that serves for packet detection, automatic gain control settling, coarse CFO and
coarse timing. The \emph{long training field} (LTF) follows: a 1.6\,$\mu$s guard interval and two
repetitions of a 3.2\,$\mu$s OFDM symbol with all 52 subcarriers loaded with known BPSK values.
Cross-correlation with the LTF gives sharp peaks for fine timing, the 64-sample repetition gives fine
CFO, and the known values give the channel estimate for equalisation. The entire synchronization and
channel estimation budget of a Wi-Fi packet is 16\,$\mu$s, and it has survived, for backward
compatibility, into 802.11n, ac, ax and be.

\mdcfig{ch10_ofdm_timing}{OFDM timing metrics. (a) The Schmidl--Cox metric for a 64-point OFDM frame
with a 16-sample cyclic prefix and a CFO of 0.37 subcarrier spacings, at two SNRs: the plateau is the
length of the CP. (b) The 802.11a/g legacy preamble at 20\,MS/s with a 100\,kHz CFO and 10\,dB SNR:
the STF's delay-and-correlate metric forms an 8\,$\mu$s plateau, and cross-correlation with the LTF
produces two sharp peaks for fine timing.}

\subsection{Cell search in LTE and 5G NR}

A phone switched on in an unknown place must find a cell, align its time and frequency to it and
learn its identity---\term{cell search}---before it can do anything else. It knows neither the carrier
frequency exactly (its crystal may be several ppm off), nor the frame timing, nor which of hundreds of
cell identities are nearby. The 5G NR solution is the \term{synchronization signal block} (SSB),
defined in 3GPP TS 38.211: four OFDM symbols by 240 subcarriers (20 resource blocks) containing
\begin{itemize}
\item the \emph{primary synchronization signal} (PSS): a length-127 BPSK m-sequence, one of three
cyclic shifts, identifying $N_{ID}^{(2)}\in\{0,1,2\}$;
\item the \emph{secondary synchronization signal} (SSS): a length-127 sequence built from two
m-sequences, one of 336 values, identifying $N_{ID}^{(1)}$; together they give the physical cell identity
$3N_{ID}^{(1)}+N_{ID}^{(2)}$, one of 1008;
\item the \emph{physical broadcast channel} (PBCH) with its demodulation reference signal, carrying
the master information block and, implicitly, the SSB index and frame timing.
\end{itemize}
Several SSBs are sent in a burst within a 5\,ms half-frame, each on a different beam (up to 4, 8 or 64
depending on frequency range), and a phone doing initial search may assume the burst repeats every
20\,ms. The search proceeds in exactly the order of this chapter: correlate against the three PSS
candidates (in time, non-coherently, over a range of frequency hypotheses) to find symbol timing,
$N_{ID}^{(2)}$ and coarse CFO; detect the SSS, now coherently, using the PSS as a phase reference, to get
the cell identity; decode the PBCH to learn the SSB index and hence the frame boundary. SSBs are placed
on a coarse \emph{synchronization raster} rather than the fine channel raster, so the phone has fewer
frequencies to try.

LTE's PSS used length-63 Zadoff--Chu sequences (roots 25, 29 and 34) and its SSS 168 interleaved
m-sequence pairs, for 504 cell identities; both are sent every 5\,ms around the centre of the carrier.
The change of PSS family in NR has a nice synchronization explanation, shown in
Figure~\ref{fig:ch10_pss_cfo}. A frequency offset of one subcarrier shifts a frequency-domain
Zadoff--Chu sequence by one position, and a shifted Zadoff--Chu sequence is the original multiplied by a
linear phase---which in the time domain is a delay. The correlation peak therefore \emph{moves} with
frequency offset: a CFO of one subcarrier spacing produces an almost full-height peak at the wrong
timing. The BPSK m-sequence of NR has no such coupling; an integer CFO simply destroys the correlation,
so a peak at the right timing can be found only at the right frequency hypothesis. With the larger
initial frequency errors expected at higher carrier frequencies, avoiding that timing--frequency
ambiguity is valuable.

\mdcfig{ch10_pss_cfo}{Time-domain correlation of the primary synchronization signal under frequency
offset (in subcarrier spacings, SCS). (a) LTE's Zadoff--Chu PSS: an offset of one subcarrier moves the
peak by about 100 samples with almost no loss in height, a false timing. (b) NR's m-sequence PSS: the
peak falls but never moves.}

\begin{table}[htbp]\centering\small
\caption{Synchronization signals of LTE and 5G NR (3GPP TS 36.211 and TS 38.211).}
\label{tab:ch10:ssb}
\begin{tabular}{lll}
\toprule
 & LTE & 5G NR \\
\midrule
PSS & Zadoff--Chu, length 63, 3 roots & BPSK m-sequence, length 127, 3 shifts \\
SSS & interleaved m-sequences, 168 values & two m-sequences, 336 values \\
Cell identities & 504 & 1008 \\
Location & 6 central RBs, every 5\,ms & SSB of 20 RBs on a sync raster \\
Beams & none & up to 4, 8 or 64 SSBs per burst \\
Frame timing from & SSS (subframe 0 or 5) & PBCH and its DMRS (SSB index) \\
\bottomrule
\end{tabular}
\end{table}

\section{Burst and continuous modems}
\label{sec:ch10:burst}

The difference between burst and continuous transmission runs through every choice in this chapter,
and it is worth stating directly.

A \emph{continuous} receiver---a DVB-S2 set-top box, a microwave backhaul link, a cable modem's
downstream---acquires once and then tracks for hours. Acquisition time is almost irrelevant; what
matters is tracking accuracy, cycle-slip rate and recovery from fades. Loops with narrow bandwidth,
long averaging and periodic unique words are the natural tools, and pilots can be sparse.

A \emph{burst} receiver---a Wi-Fi access point, an LTE base station receiving random access, a
satellite return-link hub serving thousands of terminals in time-division multiple access---must
acquire every packet from scratch, often in a few tens of symbols, and each packet may come from a
different transmitter with its own frequency, timing and power. Here overhead is the enemy: every
preamble symbol is capacity lost. Burst modems therefore favour:
\begin{itemize}
\item \emph{Feedforward, data-aided estimators} (Luise--Reggiannini frequency, ML phase,
Oerder--Meyr or data-aided timing) that converge in one pass over the preamble.
\item \emph{Store and process}: buffer the whole burst, estimate from the preamble \emph{and} the
postamble or mid-amble, and even iterate with decoder decisions (turbo synchronization). A
mid-amble, as in GSM's 26-bit training sequence in the middle of each burst, halves the maximum
extrapolation distance of the phase estimate.
\item \emph{Pre-compensation} at the transmitter, as in LTE and NR, where the base station tells each
phone how much to advance its timing (the timing advance command) and the phones lock their
oscillators to the downlink, so that uplink bursts arrive already nearly synchronized.
\end{itemize}
A software radio can do what hardware modems of the past could not: capture the entire burst,
process it as many times as needed, and pick the algorithm per packet. This is a good way to build a
first receiver with the USRP---capture, then synchronize offline---before moving to streaming.

\section{Clock and data recovery in serial links}
\label{sec:ch10:cdr}

Every high-speed wireline link---PCI Express, USB, SATA, Ethernet from 1 to 800\,Gb/s, the links
inside every data centre---sends data without a separate clock. The receiver's \term{clock and data
recovery} (CDR) circuit extracts the timing from the data transitions, a baseband version of symbol
timing recovery at speeds of 10 to 200\,GBd, where every gate delay matters.

\subsection{Bang-bang CDRs}

At these rates a linear timing error detector is hard to build, so most CDRs use a \term{bang-bang}
phase detector, often called an Alexander detector after J.\,D.\,H.~Alexander's 1975 paper. It samples
the data twice per bit: once at the centre (the data sample) and once at the expected transition (the
edge sample). Comparing the edge sample with the data samples on either side reveals, whenever a
transition occurs, whether the clock is early or late---but only the sign, not the magnitude. The
detector's output is $\pm1$ or 0, and the loop is a nonlinear, slew-rate-limited feedback system that
dithers around the lock point in steps set by the phase interpolator's resolution. Its effective gain
depends on the input jitter, which makes analysis subtle, but it is simple, fast and robust. In
modern DSP-based SerDes for PAM-4 at 50 to 100\,GBd (Chapter~\ref{ch:baseband}), the receiver samples
with an ADC at one sample per symbol, and the CDR uses a baud-rate Mueller--M\"uller detector operating
on equalised samples.

\subsection{Jitter transfer and jitter tolerance}

Serial-link standards specify the CDR through its response to sinusoidal jitter---timing modulation of
the incoming data---which is exactly the PLL's $H$ and $1-H$ (Figure~\ref{fig:ch10_cdr}(a)).
\begin{description}
\item[Jitter transfer] is $|H(f)|$: how much input jitter appears on the recovered clock. In a
repeater or retimer it must be limited, with little peaking, so that jitter does not accumulate
along a chain of regenerators---the historic SONET concern.
\item[Jitter tolerance] is the largest sinusoidal jitter at which the receiver still achieves its
target BER. Low-frequency jitter, even of many unit intervals (UI), is tracked by the CDR and does no
harm; high-frequency jitter passes through $1-H$ and consumes the eye opening directly. If the
eye's timing margin is $A$ UI, the tolerance is approximately $A/|1-H(f)|$ (Figure~\ref{fig:ch10_cdr}(b)),
rising at 20 or 40\,dB/decade below the loop bandwidth. Compliance masks are drawn on the same axes,
and the receiver passes if its tolerance curve lies above the mask.
\end{description}
The loop bandwidth is chosen to track the low-frequency wander of the transmitter's clock---including
the deliberate \emph{spread-spectrum clocking} used in PCIe and SATA to reduce electromagnetic
interference, which sweeps the clock by up to 5000\,ppm at about 30\,kHz---while filtering high-frequency
jitter. For measurement purposes, a common convention (originating in Fibre Channel) defines a
reference ``golden PLL'' with a corner frequency of the baud rate divided by 1667.

\mdcfig{ch10_cdr}{A second-order CDR loop at 10.3125\,GBd with its corner at $f_{\text{baud}}/1667$.
(a) Jitter transfer $|H|$ and error transfer $|1-H|$. (b) Jitter tolerance for a 0.3\,UI eye margin,
$0.3/|1-H(f)|$, against an illustrative compliance mask: low-frequency jitter of many UI is tracked;
high-frequency jitter must fit within the eye.}

\section{Network time and frequency}
\label{sec:ch10:network}

The synchronizers so far align one receiver with one transmitter. Modern networks also need the
\emph{transmitters} to agree with each other---on frequency, on phase and on absolute time---and that
requires distributing time and frequency across a continent.

\subsection{Why networks need synchronization}

Three requirements drive network synchronization.
\begin{itemize}
\item \textbf{Frequency.} Base stations must transmit on frequency so that phones, which lock to
them, meet their own emission limits and can hand over between cells without re-acquiring. NR
requires a base-station frequency error within $\pm0.05$\,ppm (TS 38.104). Circuit-switched networks
(T1/E1, SONET/SDH) additionally need all nodes to run at the same rate on average, or buffers overflow
and frames are slipped.
\item \textbf{Phase and time.} In time-division duplex (TDD) networks, all base stations in an area
must switch between transmit and receive at the same instant, or a transmitting base station will
blast its neighbour's receiver. TS 38.133 specifies a cell phase synchronization accuracy of 3\,$\mu$s
between base stations with overlapping coverage, which network operators budget as $\pm1.5\,\mu$s from
a common time reference---the figure in ITU-T G.8271. Positioning, carrier aggregation across sites and
coordinated multipoint transmission can require tighter alignment still.
\item \textbf{Holdover.} When the reference is lost---a GNSS antenna fault, jamming, a network
outage---the local oscillator must keep time on its own until service is restored.
\end{itemize}

\subsection{Sources: oscillators and GNSS}

The quality of an oscillator is described by its fractional frequency error and how that error
drifts, usually characterised by the Allan deviation, the standard deviation of frequency
differences between adjacent averaging intervals. Crystal oscillators range from simple XOs through
temperature-compensated (TCXO) and oven-controlled (OCXO) units; atomic references include rubidium
vapour cells, caesium beam standards (the primary references of national laboratories and the basis
of the SI second) and hydrogen masers. What matters for holdover is the accumulated time error. An
oscillator with initial fractional frequency error $y_0$ and linear drift (ageing) $D$ accumulates
\begin{equation}
x(t)=y_0\,t+\tfrac12Dt^2,
\end{equation}
plus random-walk terms. Figure~\ref{fig:ch10_holdover} shows how quickly the $\pm1.5\,\mu$s TDD budget
is consumed for a range of illustrative oscillator grades: minutes for a TCXO, hours for an OCXO, and days
for the best OCXOs and rubidium standards.

\mdcfig[0.7]{ch10_holdover}{Accumulated time error during holdover for oscillators with the stated
initial fractional frequency error $y_0$ and daily ageing $D$ (illustrative values chosen to span the
range of typical grades; real holdover also depends on temperature). The dashed line is the
$\pm1.5\,\mu$s TDD budget.}

The most common time source is a GNSS receiver. A \term{GPS-disciplined oscillator} (GPSDO) steers a
local OCXO or rubidium oscillator with a very narrow loop---time constants of minutes to hours---to the
GNSS pulse-per-second output, combining the long-term accuracy of the satellites' caesium and rubidium
clocks with the short-term stability of the local oscillator. A good GPSDO holds time to within tens of
nanoseconds of UTC. It is a PLL like every other in this chapter: the loop bandwidth is set at the
crossover between the GNSS time-transfer noise (large at short averaging times) and the oscillator's
wander (large at long times). The USRP B200 can be fitted with such a module. GNSS is, however,
vulnerable to jamming and spoofing and unavailable indoors, so networks also distribute timing over
their own links.

\subsection{Synchronous Ethernet and the Precision Time Protocol}

\term{Synchronous Ethernet} (SyncE, ITU-T G.8261/G.8262) distributes \emph{frequency} through the
physical layer: each node recovers the line clock from its upstream link with a CDR, cleans it with a
narrow PLL, and uses it to clock its downstream transmitters, exactly as SONET did. The Ethernet
equipment clock of G.8262 must hold $\pm4.6$\,ppm when free-running, and a messaging channel carries a
quality level so that nodes select the best available source.

The \term{Precision Time Protocol} (PTP, IEEE 1588) distributes \emph{time} with timestamped packets
(Figure~\ref{fig:ch10_ptp}). The master sends a Sync message at time $t_1$ (by its clock), which the
slave receives at $t_2$ (by its clock); the slave sends a Delay\_Req at $t_3$, which the master receives
at $t_4$. If the path delay $d$ is the same in both directions and the slave's clock is ahead by an
offset $o$,
\begin{gather}
t_2-t_1=d+o,\qquad t_4-t_3=d-o, \notag\\
o=\frac{(t_2-t_1)-(t_4-t_3)}{2},\qquad d=\frac{(t_2-t_1)+(t_4-t_3)}{2}.
\label{eq:ch10:ptp}
\end{gather}
The weak point is the symmetry assumption: any \emph{asymmetry} $\Delta=d_{ms}-d_{sm}$ between the two
directions causes an undetectable time error of $\Delta/2$. Packet queuing in switches produces
variable, asymmetric delay, so telecom-grade PTP uses \emph{hardware timestamping} at the physical
layer and requires every switch to participate: a \emph{boundary clock} terminates PTP and acts as a
new master, and a \emph{transparent clock} measures each packet's residence time and adds it to a
correction field. The ITU-T G.8275.1 telecom profile, with full on-path support and SyncE for
frequency, is designed to deliver the $\pm1.5\,\mu$s budget over chains of many boundary clocks; CERN's
White Rabbit extension of PTP reaches sub-nanosecond accuracy by also measuring the fibre asymmetry.

\begin{figure}[htbp]\centering
\begin{tikzpicture}[arr/.style={-{Stealth[length=2.2mm]},thick,navy},font=\small]
\draw[thick,navy] (0,0) -- (0,-4.2) node[below] {\sffamily\scriptsize master};
\draw[thick,accent] (5.2,0) -- (5.2,-4.2) node[below] {\sffamily\scriptsize slave};
\node[above,font=\sffamily\scriptsize] at (0,0) {time};
\draw[arr] (0,-0.5) node[left]{$t_1$} -- node[above,sloped,font=\scriptsize]{Sync} (5.2,-1.3) node[right]{$t_2$};
\draw[arr,densely dashed] (0,-1.0) -- node[below,sloped,font=\scriptsize]{Follow\_Up ($t_1$)} (5.2,-1.8);
\draw[arr,accent] (5.2,-2.6) node[right]{$t_3$} -- node[above,sloped,font=\scriptsize]{Delay\_Req} (0,-3.4) node[left]{$t_4$};
\draw[arr,densely dashed] (0,-3.6) -- node[below,sloped,font=\scriptsize]{Delay\_Resp ($t_4$)} (5.2,-4.1);
\end{tikzpicture}
\caption{The IEEE 1588 two-way time transfer. Four timestamps give the clock offset and the mean
path delay, provided the path is symmetric.}
\label{fig:ch10_ptp}
\end{figure}

\begin{worked}[The cost of asymmetry]
A PTP path runs over a fibre pair in which the two directions differ in length by 200\,m, perhaps because
of patching at an intermediate site. Light in fibre travels at about 5\,$\mu$s/km, so the asymmetry is
$\Delta=1\,\mu$s and the slave's time is wrong by $\Delta/2=500$\,ns, a third of the whole TDD budget,
with no way for the protocol to detect it. Asymmetry of this kind is why operators calibrate links,
prefer bidirectional single-fibre optics (where both directions share one fibre on two wavelengths), and
keep GNSS at the base station as a cross-check.
\end{worked}

\begin{inpractice}[Timing a 5G TDD network]
A typical 5G TDD network combines everything in this section. A primary reference time clock (PRTC)
at a core site, disciplined by multi-constellation GNSS and backed by a caesium or rubidium standard
for holdover, feeds PTP (G.8275.1) with SyncE over the transport network. Every router along the way is
a boundary clock with hardware timestamping. At each cell site the radio unit locks its synthesizers
to the recovered frequency and aligns its frame boundary to the recovered time, often with a local
GNSS receiver as a backup. Inside the radio, a jitter-cleaning PLL removes the packet-network wander,
the RF synthesizers of Chapter~\ref{ch:sdr} multiply the clean reference up to the carrier, and the
phones, finally, lock their own oscillators to the base station with the carrier and timing loops of
this chapter. From the caesium atom to the phone's constellation diagram, it is phase-locked loops
all the way down.
\end{inpractice}

\section{Putting it together}
\label{sec:ch10:together}

A complete single-carrier burst receiver built from the pieces of this chapter---and implemented in
the companion lab---runs the following chain on a buffer of samples at four samples per symbol:
\begin{enumerate}
\item matched filter (RRC);
\item Gardner timing recovery with a cubic Farrow interpolator, $B_LT\approx0.01$, producing one sample
per symbol;
\item coarse CFO by the fourth-power FFT estimator over the first 1024 symbols, and correction;
\item a decision-directed second-order PLL, $B_LT\approx0.01$, removing the residual frequency and
tracking phase;
\item correlation with the known preamble to find the frame start, and use of the complex correlation
peak to rotate away the loop's $\pi/2$ ambiguity;
\item detection and BER measurement.
\end{enumerate}
It works over a wide range of offsets and SNRs, and it fails in instructive ways: beyond $\pm1/8$ cycle
per symbol of CFO the fourth-power estimator aliases; below about 5\,dB the decision-directed loop
slips; and if the preamble is too short the frame synchronizer finds data instead.

\begin{keyidea}[Designing a synchronizer]
\begin{itemize}
\item List the offsets and their dynamics (Table~\ref{tab:ch10:offsets}) as fractions of the symbol
rate. They determine the architecture.
\item Order the estimators so that each tolerates the residual errors of those before it:
coarse frequency, timing, frame, fine frequency and phase.
\item Compute the CRB for your preamble and pilot budget, and check the estimators' thresholds
against the lowest SNR at which the code must work.
\item Choose loop bandwidths from $\sigma_\phi^2=B_LT/\esno$ and the dynamics to be tracked, use
type-2 loops for frequency offsets, and aid acquisition with FLLs or feedforward estimates.
\item Resolve ambiguities with known symbols, and plan for cycle slips.
\item Test with real impairments together---CFO, clock offset, phase noise, fading---not one at a
time.
\end{itemize}
\end{keyidea}

\begin{pitfall}[Testing synchronizers one impairment at a time]
A timing loop that works perfectly with no frequency offset, and a carrier loop that works perfectly
with perfect timing, may fail together, because a real receiver's loops interact: a decision-directed
timing detector makes bad decisions when the carrier loop has not converged, the carrier loop's
bandwidth changes with the timing loop's residual ISI, and both are disturbed by AGC transients at the
start of a burst. Always test the full chain with all impairments at once, with random initial
offsets, and measure acquisition failures and slip rates over many thousands of trials, not only the
steady-state jitter of a single run.
\end{pitfall}

\begin{labbox}[Synchronization from scratch]
{\raggedright Lab file: \lab{moderncomms-labs/labs/lab04\_synchronization.py}\\
Notebook: \lab{moderncomms-labs/labs/lab04\_synchronization.ipynb}\par}
\smallskip
\noindent The lab builds every synchronizer in this chapter from the \texttt{commlib.sync} module. It shows the four
impairments of Figure~\ref{fig:ch10_impairments} on a QPSK constellation; estimates CFO with the
fourth-power FFT estimator and exposes its $\pm1/8$ range; lets you vary $B_LT$, CFO, SNR and
modulation for the decision-directed PLL with interactive sliders and watch acquisition, jitter and the
$\pi/2$ ambiguity; measures the Gardner S-curve and the timing loop's convergence to a fractional
delay; detects a Zadoff--Chu preamble at 5\,dB and uses its phase to resolve the ambiguity; and finally
chains everything into the complete burst receiver of Section~\ref{sec:ch10:together}, with sliders for
CFO, delay and SNR. Its exercises ask you to implement the Luise--Reggiannini and Fitz estimators and
compare them with the MCRB (Figure~\ref{fig:ch10_freq_est}), measure acquisition time against $B_LT$,
and implement polyphase filter-bank clock recovery. With hardware, the same chain runs on samples
captured by a USRP B200 from a second radio: the two TCXOs supply real CFO and clock offset.
\end{labbox}

\section*{Summary}

\begin{itemize}
\item A coherent receiver must estimate carrier frequency, carrier phase, symbol timing and frame
position before it can detect anything. It does so with a chain of estimators, ordered so that each
tolerates the residual errors of the previous ones.
\item Maximum likelihood reduces every synchronization problem to correlating with what is known and
finding the peak. The Cram\'er--Rao bounds scale as $1/N$ for phase and timing and as $1/N^3$ for
frequency; the modified CRB is the practical yardstick. Thresholds, not bounds, usually decide what
works.
\item The phase-locked loop is the universal feedback synchronizer. Its linear model gives the noise
bandwidth $B_L$, the jitter $\sigma_\phi^2=N_0B_L/C$, and the steady-state errors by loop type; its
nonlinear behaviour gives lock-in, slow pull-in, hang-up and exponentially rare cycle slips.
\item Suppressed-carrier modulations require the carrier to be regenerated: data-aided,
decision-directed, Costas and $M$th-power detectors, at the price of squaring loss and an $M$-fold
ambiguity resolved by differential encoding, unique words or pilots.
\item Frequency estimation trades range against accuracy; coarse (FLL, $M$th power, short lag) and
fine (Luise--Reggiannini, Fitz, long lag) estimators are combined. OFDM is especially sensitive and
uses repeated preambles.
\item Timing recovery in software radios interpolates a free-running sample stream under control of a
loop driven by a Gardner, Mueller--M\"uller, early--late or ML timing error detector, or selects among
polyphase matched filters.
\item Frame synchronization correlates with sequences of good correlation---Barker, m-sequences,
Zadoff--Chu---with thresholds set by false-alarm and miss probabilities; LTE and NR cell search and the
802.11 preamble are complete synchronization systems in miniature.
\item The same loops recover clocks in SerDes links, specified by jitter transfer and tolerance, and
distribute time and frequency through networks with GNSS, SyncE and PTP to meet the $\pm1.5\,\mu$s
budget of TDD cellular systems.
\end{itemize}

\problems
\begin{enumerate}[label=\thechapter.\arabic*]
\item A 5G NR phone at 28\,GHz uses a TCXO with $\pm1$\,ppm accuracy before it locks to the network.
What is the worst-case CFO, and what fraction of the 120\,kHz subcarrier spacing is it? Can the
Schmidl--Cox half-symbol correlation resolve it unambiguously?
\item Derive the ML phase estimator~\eqref{eq:ch10:mlphase} from the likelihood~\eqref{eq:ch10:llf}
and show that its variance at high SNR is $1/(2N\esno)$, the CRB.
\item Derive the frequency CRB of Table~\ref{tab:ch10:crb} for $z_k=e^{j(2\pi\nu k+\theta)}+w_k$,
$k=0,\dots,N-1$, with $\theta$ unknown. (Hint: compute the $2\times2$ Fisher information matrix for
$(\nu,\theta)$ and invert it.) Why does the bound fall as $N^3$?
\item For the second-order loop~\eqref{eq:ch10:H2}, derive the noise bandwidth~\eqref{eq:ch10:BL} and
find the damping factor that minimises $B_L$ for a given $\omega_n$.
\item A type-2 carrier loop with $\zeta=0.707$ must track a LEO satellite downlink whose Doppler rate
reaches 3\,kHz/s while keeping the steady-state phase error below $5^\circ$. Find the minimum natural
frequency and the resulting $B_L$. What loop SNR does this imply at $C/N_0=55$\,dB-Hz, and what is the
mean time between cycle slips from~\eqref{eq:ch10:slip}?
\item Using \texttt{commlib.loop\_gains}, design a digital PLL with $B_LT=0.002$ for a 250\,ksym/s QPSK
signal. Compute $k_p$ and $k_i$, the lock-in range in hertz, and the approximate pull-in time from an
offset of 500\,Hz.
\item Show that the BPSK Costas loop's error signal is $\tfrac12\sin2\phi$ in the absence of noise, and
that the QPSK hard-limited detector $\operatorname{sgn}(I)Q-\operatorname{sgn}(Q)I$ has period $\pi/2$.
Sketch both S-curves and mark the stable and unstable lock points.
\item Show that the Luise--Reggiannini estimator is unbiased at high SNR and has range $|\nu|<1/(L+1)$.
For a 64-symbol preamble, compare its range with $L=32$ to the range of the fourth-power estimator.
Propose a two-stage estimator that covers $\pm0.1$ cycles/symbol and approaches the CRB.
\item Derive the Gardner detector's S-curve for BPSK with a raised-cosine overall pulse, and show that
its slope at $\tau=0$ is proportional to the roll-off for small $\beta$ (compare with
Figure~\ref{fig:ch10_ted_scurves}(b)).
\item A burst of 2000 symbols is received with a sample clock offset of 50\,ppm. How far does the
timing drift over the burst? Is a one-shot Oerder--Meyr estimate from the first 100 symbols sufficient
for 16-QAM with $\beta=0.35$? Use Figure~\ref{fig:ch10_timing_ber}.
\item A continuous stream uses a 24-bit sync word every 1000 bits. Design a search/verify/lock state
machine and compute (approximately) the probability of false lock and the mean time to lose lock
falsely at a channel BER of $10^{-3}$.
\item Explain, with the help of the Zadoff--Chu shift property, why an integer CFO moves the LTE PSS
correlation peak. By how many samples does it move for root $u=25$, $N=63$ and a 128-point FFT?
Compare with Figure~\ref{fig:ch10_pss_cfo}.
\item A PTP slave records $t_1=100.000\,000\,000$, $t_2=100.000\,012\,300$, $t_3=100.000\,150\,000$,
$t_4=100.000\,159\,700$ (seconds). Compute the offset and mean path delay. If the forward path is
actually 1\,$\mu$s longer than the reverse, what is the true offset?
\item \emph{(Simulation)} Extend the burst receiver of lab~4 with a Luise--Reggiannini estimator on the
preamble ahead of the PLL. Measure the probability of successful acquisition (BER below $10^{-3}$ over
the payload) against CFO and $\esno$, with and without the fourth-power estimator, and explain the
difference.
\item \emph{(Simulation)} Implement a Schmidl--Cox detector for an OFDM system with 256 subcarriers and a
CP of 18 samples in a two-path channel. Plot the timing metric, and measure the RMS timing and
frequency errors against SNR. Compare with the CRB for frequency estimation from $N/2$ repeated
samples.
\end{enumerate}

\furtherreading
\begin{itemize}
\item U.~Mengali and A.~N.~D'Andrea, \emph{Synchronization Techniques for Digital Receivers}, Plenum,
1997. The definitive treatment of ML synchronization, the MCRB and nearly every estimator in this
chapter.
\item H.~Meyr, M.~Moeneclaey and S.~A.~Fechtel, \emph{Digital Communication Receivers: Synchronization,
Channel Estimation, and Signal Processing}, Wiley, 1998. Encyclopaedic and implementation-minded.
\item F.~M.~Gardner, \emph{Phaselock Techniques}, 3rd ed., Wiley, 2005. The classic engineering text on
PLLs: acquisition, noise, cycle slips and synthesizers, written with rare clarity.
\item M.~Rice, \emph{Digital Communications: A Discrete-Time Approach}, Pearson, 2009. Chapters 7 and 8
develop discrete-time carrier and timing loops exactly as \texttt{commlib} implements them.
\item F.~M.~Gardner, ``Interpolation in digital modems---Part I: Fundamentals,'' \emph{IEEE Trans.
Commun.}, 41(3), 1993; L.~Erup, F.~M.~Gardner and R.~A.~Harris, ``Part II: Implementation and
performance,'' 41(6), 1993. The foundation of asynchronous timing recovery.
\item F.~M.~Gardner, ``A BPSK/QPSK timing-error detector for sampled receivers,'' \emph{IEEE Trans.
Commun.}, 34(5), 1986; K.~H.~Mueller and M.~M\"uller, ``Timing recovery in digital synchronous data
receivers,'' \emph{IEEE Trans. Commun.}, 24(5), 1976; M.~Oerder and H.~Meyr, ``Digital filter and square
timing recovery,'' \emph{IEEE Trans. Commun.}, 36(5), 1988.
\item M.~Luise and R.~Reggiannini, ``Carrier frequency recovery in all-digital modems for burst-mode
transmissions,'' \emph{IEEE Trans. Commun.}, 43(2/3/4), 1995; S.~Kay, ``A fast and accurate single
frequency estimator,'' \emph{IEEE Trans. ASSP}, 37(12), 1989.
\item T.~M.~Schmidl and D.~C.~Cox, ``Robust frequency and timing synchronization for OFDM,''
\emph{IEEE Trans. Commun.}, 45(12), 1997. Short, readable and still the starting point for OFDM
synchronization.
\item J.~L.~Massey, ``Optimum frame synchronization,'' \emph{IEEE Trans. Commun.}, 20(2), 1972.
\item f.~j.~harris and M.~Rice, ``Multirate digital filters for symbol timing synchronization in
software defined radios,'' \emph{IEEE J. Sel. Areas Commun.}, 19(12), 2001. The polyphase clock sync
used in GNU Radio.
\item A.~J.~Viterbi, \emph{Principles of Coherent Communication}, McGraw-Hill, 1966. The PLL in noise,
including the cycle-slip analysis, from the engineer who solved it.
\item 3GPP TS 38.211 (NR physical channels, PSS/SSS/SSB); ETSI EN 302 307-1 (DVB-S2 framing and pilots);
IEEE Std 1588-2019 (PTP); ITU-T G.8262 (SyncE), G.8271 and G.8275.1 (time synchronization for
telecom networks).
\end{itemize}
